\documentclass[a4paper,11pt]{article} % Reduced to 11pt for significant space saving

% ==========================================
% CORE PACKAGES & GEOMETRY
\usepackage[margin=9mm, top=4mm, bottom=4mm, includeheadfoot, headheight=4pt]{geometry}
\usepackage[T1]{fontenc}
\usepackage[utf8]{inputenc}
\usepackage[english]{babel}

% ==========================================
% MATH & FONTS
% ==========================================
\usepackage{mathtools,amssymb} % mathtools loads amsmath internally
\allowdisplaybreaks % Allows multi-line derivations to break across pages

% ==========================================
% COLORS & GRAPHICS (B&W Optimized)
% ==========================================
\usepackage{xcolor,microtype}
% Redefined colors to grayscale for B&W printing
\definecolor{primaryblue}{gray}{0.2}
\definecolor{warningcolor}{gray}{0.4}
\definecolor{errorcolor}{gray}{0.3}
\definecolor{infocolor}{gray}{0.2}
\definecolor{notecolor}{gray}{0.5}

% ==========================================
% SECTION & PAGE FORMATTING
% ==========================================
\usepackage{titlesec,fancyhdr,enumitem}

% Default heading spacing
\titleformat{\section}{\bfseries\Large}{\thesection.}{0.4em}{}
\titleformat{\subsection}{\bfseries\large}{\thesubsection.}{0.4em}{}
\titleformat{\subsubsection}{\bfseries\normalsize}{\thesubsubsection.}{0.4em}{}

\pagestyle{fancy}
\fancyhead{}
\fancyhead[R]{\leftmark}
\renewcommand{\headrulewidth}{0.1pt}
\setlength{\headsep}{10pt}

% Aggressive list compression
\setlist[itemize]{label=\textbullet,leftmargin=1.5em,nosep, topsep=2pt, partopsep=0pt}
\setlist[enumerate]{leftmargin=1.5em,labelsep=0.5em,nosep, topsep=2pt, partopsep=0pt}

% ==========================================
% COLOR BOXES (Optimized Padding & B&W)
% ==========================================
\usepackage{tcolorbox}
\tcbuselibrary{breakable,skins}

% Reduced internal padding (left/right to 3mm, top/bottom to 2mm)
\newtcolorbox{custombox}[3][]{
  enhanced, breakable, rounded corners, left=3mm, right=3mm, top=2mm, bottom=2mm,
  colframe=#2, colback=white, fonttitle=\bfseries,
  title={#3}, attach boxed title to top center={yshift=-2mm},
  boxed title style={size=small, colback=#2, colframe=#2, sharp corners}, #1
}

\newenvironment{important}{\begin{custombox}{errorcolor}{Important}}{\end{custombox}}
\newenvironment{notebox}[1][]{\begin{custombox}[#1]{notecolor}{Note}}{\end{custombox}}
\newenvironment{warningbox}[1][]{\begin{custombox}[#1]{warningcolor}{Warning}}{\end{custombox}}
\newenvironment{summary}{\begin{custombox}{infocolor}{Summary}}{\end{custombox}}
\newenvironment{summarybox}[1][]{\begin{custombox}[#1]{infocolor}{Final answer}}{\end{custombox}}

% Compact problem statement box
\newtcolorbox{problembox}[1][]{
  enhanced, breakable, left=3mm, right=3mm, top=1.5mm, bottom=1.5mm,
  colframe=black, colback=gray!5, fonttitle=\bfseries,
  title={Problem}, boxrule=0.5pt, #1
}

% ==========================================
% HYPERLINKS & REFERENCES
% ==========================================
\usepackage[hidelinks]{hyperref}
\usepackage{cleveref}

% ==========================================
% GLOBAL MATH SPACING COMPRESSION
% ==========================================
\AtBeginDocument{
  \setlength{\abovedisplayskip}{4pt plus 2pt minus 2pt}
  \setlength{\belowdisplayskip}{4pt plus 2pt minus 2pt}
  \setlength{\abovedisplayshortskip}{2pt plus 1pt minus 1pt}
  \setlength{\belowdisplayshortskip}{2pt plus 1pt minus 1pt}
}
\setlength{\parskip}{3pt plus 1pt minus 1pt} % Added slight paragraph spacing for readability
\setlength{\parindent}{0pt} % Removed indent to save horizontal space


\begin{document}

\makeatletter
\@starttoc{toc}
\makeatother


% ---------------------------------------------------------------------
% PROBLEM 1
% ---------------------------------------------------------------------
\section{First-Order PDEs and Classification}

\begin{notebox}[title={General Procedure: Method of Characteristics}]
First-order PDEs of the form $P(x,y,u)\,u_x + Q(x,y,u)\,u_y = R(x,y,u)$ are solved using the \textbf{method of characteristics}, which reduces the PDE to a system of ODEs along characteristic curves.
\begin{enumerate}
    \item \textbf{Write the Lagrange--Charpit (auxiliary) equations:} $\displaystyle\frac{dx}{P} = \frac{dy}{Q} = \frac{du}{R}.$
    \item \textbf{Extract two independent invariants.} Pair ratios to form integrable ODEs. Integrate each to obtain constants $C_1 = \phi(x,y)$ and $C_2 = \psi(x,y,u)$.
    \item \textbf{Form the general solution.} Set $C_2 = f(C_1)$, i.e.\ $\psi = f(\phi)$, where $f$ is an arbitrary function.
    \item \textbf{Apply the initial/boundary condition.} Substitute the given curve data to determine $f$ explicitly, then replace its argument to obtain the particular solution $u(x,y)$.
    \item \textbf{Verify.} Check that the solution satisfies both the PDE and the initial/boundary curve.
\end{enumerate}
\textbf{For classification of 2nd-order PDEs} $Au_{xx}+2Bu_{xy}+Cu_{yy}+\cdots=0$:
\begin{enumerate}
    \item Compute the discriminant $\Delta = B^2-AC$. If $\Delta>0$: Hyperbolic; $\Delta=0$: Parabolic; $\Delta<0$: Elliptic.
    \item Find characteristic curves by solving $A\lambda^2 - 2B\lambda + C = 0$ where $\lambda = dy/dx$; integrate to get new coordinates $(\xi,\eta)$.
    \item Transform the PDE to $(\xi,\eta)$ via chain rule; the equation reduces to its canonical form (e.g.\ $u_{\xi\eta}=0$ for hyperbolic).
    \item Integrate the canonical form and apply Cauchy data to determine the arbitrary functions.
\end{enumerate}
\end{notebox}

\subsection{\texorpdfstring{$x u_x + y u_y = 2u$; $u(x,1)=x^2+x$}{Solve x u_x + y u_y = 2u by characteristics; particular solution satisfying u(x,1)=x^2+x}}

\begin{problembox}
Consider the quasilinear equation:
$$x \frac{\partial u}{\partial x} + y \frac{\partial u}{\partial y} = 2u$$

(a) Find the general solution $u(x, y)$ using the method of characteristics.

(b) Find the particular solution satisfying the curve:
$$u(x, 1) = x^2 + x$$
\end{problembox}


\textbf{(a) General Solution.}
The auxiliary equations $\frac{dx}{x} = \frac{dy}{y} = \frac{du}{2u}$ give two invariants:
\begin{itemize}
    \item $\frac{dx}{x}=\frac{dy}{y} \Rightarrow C_1 = y/x$
    \item $\frac{dx}{x}=\frac{du}{2u} \Rightarrow C_2 = u/x^2$
\end{itemize}
Setting $C_2=f(C_1)$: $\displaystyle\boxed{u(x,y)=x^2 f\!\left(\tfrac{y}{x}\right)}$.

\textbf{(b) Particular Solution}

We apply the condition $u(x, 1) = x^2 + x$. Substituting $y=1$ and $u = x^2+x$ into the general solution:
$$ x^2 + x = x^2 f\left(\frac{1}{x}\right) $$
Dividing by $x^2$ (assuming $x \neq 0$):
$$ 1 + \frac{1}{x} = f\left(\frac{1}{x}\right) $$
To determine the structure of $f$, let $t = \frac{1}{x}$, giving $f(t) = 1 + t$.
Substituting $t = \frac{y}{x}$ back: $u(x, y) = x^2 + xy$.

\newpage
\subsection{\texorpdfstring{$xz\,z_x + yz\,z_y = xe^{x+y}$; $z^2=2e^{x+1}$ on $y=1$}{Solve quasilinear PDE xz z_x + yz z_y = xe^{x+y} with BC z^2=2e^{x+1} at y=1}}


\begin{problembox}

   $$ xz \frac{\partial z}{\partial x} + yz \frac{\partial z}{\partial y} = x e^{x+y}$$
 boundary condition:
    $z^2 = 2e^{x+1} \quad \text{on the curve} \quad y = 1$
\end{problembox}
With $P=xz$, $Q=yz$, $R=xe^{x+y}$, the Lagrange equations are $\frac{dx}{xz}=\frac{dy}{yz}=\frac{dz}{xe^{x+y}}$.

\textbf{First invariant:} $\frac{dx}{xz}=\frac{dy}{yz}\Rightarrow \frac{dx}{x}=\frac{dy}{y}\Rightarrow u=\tfrac{x}{y}=c_1$.

\textbf{Second Characteristic ($v$):} From the first and third ratios, $e^{x+y}\,dx = z\,dz$.
Using $c_1=x/y$ to write $y=x/c_1$:
\begin{equation}
    e^{x + \frac{x}{c_1}} dx = z dz \implies e^{x\left(1 + \frac{1}{c_1}\right)} dx = z dz
\end{equation}
Integrating:
\begin{equation}
    \int e^{x\left(\frac{c_1+1}{c_1}\right)} dx = \int z dz
\end{equation}
\begin{equation}
    \frac{c_1}{c_1+1} e^{x\left(\frac{c_1+1}{c_1}\right)} = \frac{1}{2}z^2 + C_2
\end{equation}
Substituting $c_1 = x/y$:
\begin{equation}
    \frac{x/y}{(x/y)+1} = \frac{x}{x+y}, \quad x\left(\frac{x/y+1}{x/y}\right) = x+y
\end{equation}
The equation becomes:
\begin{equation}
    \frac{x}{x+y} e^{x+y} - \frac{1}{2}z^2 = -C_2
\end{equation}
The general solution $v=f(u)$ is:
\begin{equation} \label{eq:general_sol}
    z^2 = \frac{2x e^{x+y}}{x+y} + f\left(\frac{x}{y}\right)
\end{equation}

\textbf{Part B: Particular Solution.}
Apply B.C.\ $y=1$, $z^2=2e^{x+1}$ in \eqref{eq:general_sol}:
\begin{equation}
    2e^{x+1} = \frac{2x e^{x+1}}{x+1} + f\left(\frac{x}{1}\right)
\end{equation}
Solving for $f$:
\begin{equation}
    f(\xi) = \frac{2e^{\xi+1}}{\xi+1}
\end{equation}
With $\xi=x/y$:
\begin{equation}
    f\left(\frac{x}{y}\right) = \frac{2y e^{\frac{x+y}{y}}}{x+y}
\end{equation}
Substituting into \eqref{eq:general_sol}:
\begin{equation}
    z^2 = \frac{2x e^{x+y}}{x+y} + \frac{2y e^{\frac{x+y}{y}}}{x+y}
\end{equation}

\textbf{Final Result:}
The unique integral surface through $\Gamma$:
\begin{equation}
    \boxed{ z^2 = \frac{2}{x+y} \left( x e^{x+y} + y e^{\frac{x}{y} + 1} \right) }
\end{equation}

\textbf{Checks.} At $y=1$:
\begin{align*}
    z^2 \big|_{y=1} &= \frac{2}{x+1} \left( x e^{x+1} + 1 \cdot e^{x + 1} \right) \\
    &= \frac{2}{x+1} e^{x+1} (x+1) \\
    &= 2e^{x+1} \checkmark
\end{align*}
\textbf{Singularity:} undefined on $x=-y$.
% ---------------------------------------------------------------------
% PROBLEM 2
% ---------------------------------------------------------------------
\newpage
\subsection{\texorpdfstring{$u_{xx} - 4u_{xy} + 3u_{yy} = 0$; classify, reduce to canonical form, solve with $u(x,0)=\sin x$, $u_y(x,0)=\cos x$}{Classify hyperbolic PDE, find characteristic coords, solve Cauchy problem}}

\begin{problembox}
(a) Classify the equation and reduce it to its canonical form:
$$u_{xx} - 4u_{xy} + 3u_{yy} = 0$$

(b) Find the general solution $u(x,y)$.

(c) Determine the solution given: $u(x, 0) = \sin x, \quad u_y(x, 0) = \cos x$.
\end{problembox}


\textbf{(a) Classification and Reduction}

\textbf{Step 1: Classify the equation.}
The discriminant is $\Delta = B^2 - 4AC = (-4)^2 - 4(1)(3) = 16 - 12 = 4$. Since $\Delta > 0$, the equation is \textbf{Hyperbolic}.

\textbf{Step 2: Find characteristic curves.}
The characteristic slopes $\lambda = \frac{dy}{dx}$ satisfy $\lambda^2 - 4\lambda + 3 = 0 \implies (\lambda-3)(\lambda-1)=0$.
The characteristics are $\frac{dy}{dx} = 3$ and $\frac{dy}{dx} = 1$.
Integrating yields the characteristic families: $y - 3x = c_1$ and $y - x = c_2$.

\textbf{Step 3: Characteristic coordinates} $\xi = 3x+y$, $\eta = x+y$:
$$u_x = 3u_\xi + u_\eta, \quad u_y = u_\xi + u_\eta$$
$$u_{xx} = 9u_{\xi\xi} + 6u_{\xi\eta} + u_{\eta\eta}$$
$$u_{xy} = 3u_{\xi\xi} + 4u_{\xi\eta} + u_{\eta\eta}$$
$$u_{yy} = u_{\xi\xi} + 2u_{\xi\eta} + u_{\eta\eta}$$

\textbf{Step 4: Substitute:}
$$ (9 - 12 + 3)u_{\xi\xi} + (6 - 16 + 6)u_{\xi\eta} + (1 - 4 + 3)u_{\eta\eta} = 0 $$
$$ 0 \cdot u_{\xi\xi} - 4u_{\xi\eta} + 0 \cdot u_{\eta\eta} = 0 \implies u_{\xi\eta} = 0 $$
Canonical form: $u_{\xi\eta} = 0$.

\textbf{(b) General Solution.}
Integrating $u_{\xi\eta}=0$: $u(x,y)=F(3x+y)+G(x+y)$.

\textbf{(c) Specific Solution}
Using initial data at $y=0$:
1. $u(x,0) = F(3x) + G(x) = \sin x$
2. $u_y(x,0) = F'(3x)\cdot 1 + G'(x)\cdot 1 = \cos x$

Differentiating (1) w.r.t $x$:
3. $3F'(3x) + G'(x) = \cos x$

Subtracting (2) from (3):
$$ 2F'(3x) = 0 \implies F(3x) = C_1 \text{ (constant)} $$
Since $F$ is constant, from (1): $C_1 + G(x) = \sin x \implies G(x) = \sin x - C_1$.
Substituting back into the general solution:
$$ u(x,y) = \sin(x+y) $$

\newpage
% ---------------------------------------------------------------------
% NEW PROBLEM (Section 1)
% ---------------------------------------------------------------------
\subsection{\texorpdfstring{$u_{xx} + 2u_{xy} + u_{yy} = 0$; classify, reduce to canonical form, solve with $u(x,0)=\sin x$, $u_y(x,0)=0$}{Classify parabolic PDE; single repeated characteristic; canonical form $u_{\xi\xi}=0$; general solution $xF(y-x)+G(y-x)$}}

\begin{problembox}
(a) Classify the PDE and reduce it to its canonical form:
$$u_{xx} + 2u_{xy} + u_{yy} = 0$$

(b) Find the general solution $u(x,y)$.

(c) Determine the particular solution satisfying the Cauchy data:
$$u(x, 0) = \sin x, \quad u_y(x, 0) = 0.$$
\end{problembox}

\textbf{(a) Classification and Canonical Form.} $A=1$, $B=2$, $C=1$: $\Delta = B^2-4AC = 0$, so the equation is \textbf{Parabolic}. The characteristic ODE $(dy-dx)^2=0$ gives $dy/dx=1$, i.e.\ the single repeated family $y-x=c$.

\textbf{(b) General Solution.}
Integrating twice: $u(\xi,\eta)=\xi F(\eta)+G(\eta)$, i.e.\
$$ u(x,y) = x\,F(y-x) + G(y-x). $$

\textbf{(c) Particular Solution.}
(1) $u(x,0)=xF(-x)+G(-x)=\sin x$; \quad (2) $u_y(x,0)=xF'(-x)+G'(-x)=0$.
Differentiating (1): $F(-x)-xF'(-x)-G'(-x)=\cos x$. Adding to (2): $F(-x)=\cos x$, so $F(s)=\cos s$.
Then $G(-x)=\sin x-x\cos x$, so $G(s)=-\sin s+s\cos s$.
\begin{align*}
u(x,y)&=x\cos(y-x)+(y-x)\cos(y-x)-\sin(y-x)=y\cos(y-x)-\sin(y-x).
\end{align*}

\begin{summarybox}
$u(x,y)=y\cos(y-x)-\sin(y-x)$
\end{summarybox}

\textbf{Verification:}
\begin{itemize}
    \item $u(x,0)=-\sin(-x)=\sin x$. \checkmark
    \item $u_y(x,0)=\cos(-x)+0\cdot(-\sin(-x))-\cos(-x)=0$. \checkmark
\end{itemize}

\newpage

% ---------------------------------------------------------------------
% PROBLEM 3
% ---------------------------------------------------------------------
\section{Heat Equation}
\vspace{-1cm}
\begin{notebox}[title={General Procedure: Solving the Heat Equation}]
The heat equation $u_t = k\,u_{xx}$ (possibly with lower-order terms or source terms) is solved differently depending on the domain.

\textbf{Bounded Domain --- Separation of Variables:}
\begin{enumerate}
    \item \textbf{Steady-state decomposition.} Write $u = v(x) + w(x,t)$, where $v$ satisfies $v_t=0$ with the non-homogeneous BCs and/or source.
    \item \textbf{Solve for the steady state $v(x)$:} an ODE boundary-value problem.
    \item \textbf{Formulate the transient problem} for $w$: now with homogeneous BCs and modified IC $w(x,0)=u(x,0)-v(x)$.
    \item \textbf{Separate variables} $w = X(x)T(t)$. Solve the spatial eigenvalue problem $X''+\lambda X = 0$ with BCs to obtain eigenfunctions $X_n$ and eigenvalues $\lambda_n$.
    \item \textbf{Solve the temporal ODE} $T_n' + k\lambda_n T_n = q_n(t)$ for each mode (with forcing if a source is present).
    \item \textbf{Expand the IC} in eigenfunctions to determine Fourier coefficients. Assemble $u = v + \sum_n T_n(t)X_n(x)$.
\end{enumerate}
\textbf{Semi-Infinite / Infinite Domain:}
\begin{enumerate}
    \item Reduce to a homogeneous IC via substitution (e.g.\ $\theta = u - T_0$).
    \item Apply a \textbf{Laplace transform} in $t$ (or Fourier transform in $x$) to reduce the PDE to an ODE.
    \item Solve the ODE, enforce boundedness/decay at infinity.
    \item Invert the transform using known pairs (e.g.\ $\operatorname{erfc}$ for the Laplace approach).
\end{enumerate}
\textbf{Useful substitutions:} Advection term $v_d u_x$: use Galilean frame $\xi=x-v_d t$.  Reaction term $bu$: set $u=e^{bt}w$.
\end{notebox}

\subsection{\texorpdfstring{$u_t = u_{xx} - 2u + x$ on $0<x<\pi$; $u(0,t)=0$, $u(\pi,t)=\tfrac{\pi}{2}$, $u(x,0)=\tfrac{x}{2}+\sin(3x)$}{Solve heat-reaction PDE via steady-state decomposition; exponential decay transient}}

\begin{problembox}
Solve on $0 < x < \pi$:
$$\frac{\partial u}{\partial t} = \frac{\partial^2 u}{\partial x^2} - 2u + x$$
BCs: $u(0, t) = 0, \quad u(\pi, t) = \frac{\pi}{2}$.
ICs: $u(x, 0) = \frac{x}{2} + \sin(3x)$.
\end{problembox}


Decompose $u=v(x)+w(x,t)$ where $v$ is the steady state.

\textbf{Step 1: Steady state.} $v''-2v+x=0$, $v(0)=0$, $v(\pi)=\pi/2$.
Try $v_p=Ax+B$: $-2(Ax+B)+x=0\Rightarrow A=1/2$, $B=0$. Then $v_p=x/2$ satisfies both BCs, so $v(x)=x/2$.

\textbf{Step 2: Transient.} Substituting $u=v+w$: $w_t=w_{xx}-2w$ with $w(0,t)=w(\pi,t)=0$ and $w(x,0)=\sin(3x)$.

\textbf{Step 3: Separation.} $w=X(x)T(t)$: $X''+(\lambda-2)X=0$, $X(0)=X(\pi)=0$. Eigenfunctions $\sin(nx)$ with $\lambda_n=n^2+2$, so $T_n(t)=e^{-(n^2+2)t}$.
From IC $w(x,0)=\sin(3x)$: only $n=3$, $b_3=1$.

\begin{summarybox}
$$ u(x,t) = \frac{x}{2} + e^{-11t}\sin(3x) $$
\end{summarybox}
\hrule
\subsection{\texorpdfstring{(A) Solve $n_t + v_d n_x = a^2 n_{xx}$ (advection-diffusion); (B) solve $n_t = a^2 n_{xx} + bn$ (reaction-diffusion); $n(x,0)=f(x)$}{(A) Galilean transform removes drift; (B) exponential substitution removes reaction term}}

We derive both solutions using $G(x,y;t)=\frac{1}{2a\sqrt{\pi t}}e^{-(x-y)^2/(4a^2t)}$.

\textbf{(A) Advection-Diffusion.} Use the Galilean frame $\xi=x-v_d t$, $\tau=t$. Then $\partial_t\to\partial_\tau-v_d\partial_\xi$ and the equation reduces to $n_\tau=a^2 n_{\xi\xi}$. Solving and transforming back:
\begin{equation}
    \boxed{ n(x,t) = \int_{-\infty}^{\infty} G(x-v_d t,\,y;\,t)\,f(y)\,dy }
\end{equation}

\textbf{(B) Reaction-Diffusion.} Substitute $n=e^{bt}u$. Then $n_t=e^{bt}(u_t+bu)$ and $n_{xx}=e^{bt}u_{xx}$, so the $bn$ terms cancel giving $u_t=a^2u_{xx}$ with $u(x,0)=f(x)$. Hence:
\begin{equation}
    \boxed{ n(x,t) = e^{bt}\int_{-\infty}^{\infty} G(x,y;t)\,f(y)\,dy }
\end{equation}


\newpage
\subsection{\texorpdfstring{$n_t = a^2 n_{xx} + \sin(t)\cos^2(x)$ on $0<x<\pi$; $n_x(0,t)=n_x(\pi,t)=0$, $n(x,0)=\cos x$}{Heat eq with oscillating source; Neumann BCs; three active modes k=0,1,2}}

\begin{problembox}
    $$\frac{\partial n}{\partial t} = a^2 \frac{\partial^2 n}{\partial x^2} + \sin(t)\cos^2(x),\quad \left. \frac{\partial n}{\partial x} \right|_{x=0} = \left. \frac{\partial n}{\partial x} \right|_{x=\pi} = 0,\quad n(x, t=0) = \cos(x)$$
\end{problembox}

\textbf{Step 1: Find the eigenfunctions.}
The spatial part of the homogeneous equation with Neumann BCs gives:
$$X''(x) + \lambda X(x) = 0,\quad X'(0) = X'(\pi) = 0 \implies X_k(x) = \cos(kx), \quad k = 0, 1, 2, \dots$$
with eigenvalues $\lambda_k = k^2$.

\textbf{Step 2: Expand in eigenfunctions.}
We seek a solution of the form:
$$n(x,t) = \sum_{k=0}^{\infty} T_k(t) \cos(kx)$$

\textbf{Step 3: Expand the source term.}
Using the identity $\cos^2(x) = \tfrac{1}{2} + \tfrac{1}{2}\cos(2x)$:
\[S(x,t) = \sin(t)\cos^2(x) = \tfrac{1}{2}\sin(t) + \tfrac{1}{2}\sin(t)\cos(2x)\]
So the source excites only the $k=0$ and $k=2$ modes.

\textbf{Step 4: Determine initial conditions for each mode.}
From $n(x,0) = \cos(x)$, only the $k=1$ mode is initially present: $T_1(0) = 1$ and $T_k(0) = 0$ for all $k \neq 1$.

\textbf{Step 5: Solve the ODEs.}
Substituting the series into the PDE yields, for each $k$:
\[T_k'(t) + a^2 k^2 T_k(t) = q_k(t)\]
Case 1: \(k=0\)

ODE: \(T_0'(t) = \tfrac12\sin t\), IC: \(T_0(0) = 0\).

\[
T_0(t) = \int_0^t \tfrac12 \sin\tau\,d\tau 
= \tfrac12[-\cos\tau]_0^t 
= \tfrac12(1-\cos t).
\]

Case 2: \(k=1\)

ODE: \(T_1'(t) + a^2 T_1(t) = 0\), IC: \(T_1(0) = 1\).

\[
T_1(t) = e^{-a^2 t}.
\]

Case 3: \(k=2\)

ODE: \(T_2'(t) + 4a^2 T_2(t) = \tfrac12\sin t\), IC: \(T_2(0) = 0\).

General solution: \(T_2 = C e^{-4a^2 t} + A\sin t + B\cos t\).

Solving for \(A,B,C\) with \(T_2(0)=0\) gives
\[
T_2(t) = \frac{e^{-4a^2 t} + 4a^2\sin t - \cos t}{2(16a^4+1)}.
\]

Case 4: \(k>2\)

ODE: \(T_k' + a^2 k^2 T_k = 0\), IC: \(T_k(0)=0\) \(\Rightarrow T_k(t)=0\).

Final solution (using nonzero modes \(k=0,1,2\)):

\begin{summarybox}
$  n(x,t) = \frac12(1-\cos t)
+ e^{-a^2 t}\cos x
+ \frac{e^{-4a^2 t} + 4a^2\sin t - \cos t}{2(16a^4+1)}\cos(2x).
$
  
\end{summarybox}


\newpage
\subsection{\texorpdfstring{Solve $u_t = k\nabla^2 u$ in sphere $0<r<R$; $u(R,t)=0$, $u(r,0)=f(r)$}{Spherical heat equation: substitute v=ru to reduce to 1D heat eq; Fourier-sine series coefficients}}

\begin{problembox}
Solve the heat equation in a sphere of radius $R$:
$$u_t = k\nabla^2 u, \quad 0<r<R, \quad u(R,t)=0, \quad u(r,0)=f(r)$$
\end{problembox}

\textbf{Step 1: Substitution to reduce dimension.}
Let $v(r,t)=ru(r,t)$; then
\[
    \nabla^2 u = \frac{1}{r}v_{rr}, \quad \text{so } v_t = k\,v_{rr}, \quad 0<r<R.
\]

\textbf{Conditions for $v$:} $v(0,t)=v(R,t)=0$ (finiteness/BC), $v(r,0)=rf(r)$.

\textbf{Separation of Variables.} With $v=\Phi(r)T(t)$:
\begin{equation}
    \frac{T'(t)}{k T(t)} = \frac{\Phi''(r)}{\Phi(r)} = -\lambda^2
\end{equation}
This gives $T(t)\propto e^{-k\lambda^2 t}$ and $\Phi''+\lambda^2\Phi=0$ with $\Phi(0)=\Phi(R)=0$, so $\Phi(r)=\sin(\lambda r)$ and $\lambda_n=n\pi/R$.

\begin{summarybox}[title={General Solution}]
Superposing and using $u=v/r$:
$$ u(r,t) = \frac{1}{r} \sum_{n=1}^{\infty} A_n \sin\!\left(\frac{n\pi r}{R}\right) e^{-k \left(\frac{n\pi}{R}\right)^2 t} $$
where the Fourier sine coefficients of $rf(r)$ are:
$$ A_n = \frac{2}{R} \int_0^R r f(r) \sin\left(\frac{n\pi r}{R}\right) \, dr $$
\end{summarybox}

\newpage
% ---------------------------------------------------------------------
% NEW PROBLEM (Section 2)
% ---------------------------------------------------------------------
\subsection{\texorpdfstring{$u_t = k\,u_{xx}$ for $x>0$; $u(0,t)=T_s$, $u(x,0)=T_0$, $u\to T_0$ as $x\to\infty$}{Semi-infinite heat equation: Laplace transform in time reduces to ODE; inverse gives error-function solution}}

\begin{problembox}
Solve the heat equation on the semi-infinite domain $x > 0$:
$$\frac{\partial u}{\partial t} = k\,\frac{\partial^2 u}{\partial x^2}, \quad x > 0,\; t > 0$$
subject to:
\begin{itemize}
    \item Boundary condition: $u(0, t) = T_s$ (constant surface temperature)
    \item Initial condition: $u(x, 0) = T_0$ (uniform initial temperature)
    \item Far-field condition: $u(x, t) \to T_0$ as $x \to \infty$
\end{itemize}
\end{problembox}

\textbf{Step 1: Reduce to a homogeneous IC.}
Let $\theta(x,t) = u(x,t) - T_0$. Then $\theta$ satisfies:
\[
    \theta_t = k\,\theta_{xx}, \quad \theta(0,t) = \Delta T \equiv T_s - T_0, \quad \theta(x,0)=0, \quad \theta\to 0 \text{ as } x\to\infty.
\]

\textbf{Step 2: Laplace transform.} Taking $\hat\theta=\mathcal{L}\{\theta\}$ with $\theta(x,0)=0$:
\[
    s\hat\theta = k\hat\theta_{xx} \implies \hat\theta_{xx} - \tfrac{s}{k}\hat\theta = 0.
\]
General solution: $\hat\theta = A(s)e^{-\sqrt{s/k}\,x} + B(s)e^{+\sqrt{s/k}\,x}$. The far-field condition forces $B=0$, and $\hat\theta(0,s)=\Delta T/s$ gives $A=\Delta T/s$, so
\[
    \hat\theta(x,s) = \frac{\Delta T}{s}\,e^{-\frac{x}{\sqrt{k}}\sqrt{s}}.
\]

\textbf{Step 3: Invert via a known transform pair.}
We use the Laplace transform pair:
\[
    \mathcal{L}^{-1}\!\left\{\frac{1}{s}\,e^{-a\sqrt{s}}\right\} = \mathrm{erfc}\!\left(\frac{a}{2\sqrt{t}}\right),
\]
with $a = x/\sqrt{k}$. This gives:
\[
    \theta(x,t) = \Delta T\,\mathrm{erfc}\!\left(\frac{x}{2\sqrt{kt}}\right).
\]

\textbf{Step 4: Return to original variable.}

\begin{summarybox}
$$u(x,t) = T_0 + (T_s - T_0)\,\mathrm{erfc}\!\left(\frac{x}{2\sqrt{kt}}\right)$$
\end{summarybox}

where $\mathrm{erfc}(\zeta) = \frac{2}{\sqrt{\pi}}\int_\zeta^\infty e^{-\eta^2}\,d\eta$ is the complementary error function.

\textbf{Checks:}
\begin{itemize}
    \item At $x=0$: $\mathrm{erfc}(0) = 1$, so $u(0,t) = T_0 + (T_s-T_0)\cdot 1 = T_s$. \checkmark
    \item At $t=0^+$: for $x>0$, the argument $\to\infty$ so $\mathrm{erfc}(\infty)=0$, hence $u(x,0^+)=T_0$. \checkmark
    \item As $x\to\infty$: $u \to T_0$. \checkmark
\end{itemize}

% PROBLEM 4
% ---------------------------------------------------------------------
\newpage
\section{Wave Equation}

\begin{notebox}[title={General Procedure: Solving the Wave Equation}]
The wave equation $u_{tt}=c^2 u_{xx}$ (with possible forcing) is approached according to the domain type.

\textbf{Bounded Domain --- Separation of Variables:}
\begin{enumerate}
    \item \textbf{Ansatz} $u = X(x)T(t)$ with separation constant $-\lambda$.
    \item \textbf{Spatial eigenvalue problem:} Solve $X''+\lambda X = 0$ with BCs to get eigenvalues $\lambda_n = (n\pi/L)^2$ and eigenfunctions $X_n = \sin(n\pi x/L)$.
    \item \textbf{Temporal ODE:} $T_n'' + \omega_n^2 T_n = 0$ with $\omega_n = n\pi c/L$ gives $T_n = A_n\cos(\omega_n t)+B_n\sin(\omega_n t)$.
    \item \textbf{Superpose:} $u = \sum_n [A_n\cos(\omega_n t)+B_n\sin(\omega_n t)]\sin(n\pi x/L)$.
    \item \textbf{Determine $A_n$, $B_n$} from ICs $u(x,0)=f(x)$ and $u_t(x,0)=g(x)$ via Fourier sine integrals.
\end{enumerate}
\textbf{Forced Wave Equation:}
\begin{enumerate}
    \item Expand $u$ and the source in eigenfunctions; project by orthogonality to get decoupled ODEs for each mode $T_n$.
    \item \textbf{Non-resonant} ($\Omega\ne\omega_n$): bounded beats solution.
    \item \textbf{Resonant} ($\Omega=\omega_n$): secular growth $\propto t\sin(\omega_n t)$.
\end{enumerate}
\textbf{Unbounded Domain --- d'Alembert's Formula:}
$$u(x,t) = \frac{\phi(x+ct)+\phi(x-ct)}{2}+\frac{1}{2c}\int_{x-ct}^{x+ct}\psi(s)\,ds$$
Derived by transforming to characteristic coordinates $\xi=x+ct$, $\eta=x-ct$, which reduces the PDE to $u_{\xi\eta}=0$.
\end{notebox}

\subsection{\texorpdfstring{ $u_{tt}=c^2 u_{xx}$ on $(0,L)$; $u(0,t)=u(L,t)=0$; $u(x,0)=f(x)$, $u_t(x,0)=g(x)$}{Reference: 1D wave equation IBVP via separation of variables; Fourier sine series}}

\begin{problembox}
Consider a homogeneous string of length $L$ with fixed ends:
\begin{align*}
    \text{PDE:} & \quad \frac{\partial^2 u}{\partial t^2} = c^2 \frac{\partial^2 u}{\partial x^2}, \quad x \in (0, L), \ t > 0 \\
    \text{BCs:} & \quad u(0,t) = 0, \quad u(L,t) = 0, \quad \forall t \geq 0 \\
    \text{ICs:} & \quad u(x,0) = f(x), \quad \frac{\partial u}{\partial t}(x,0) = g(x)
\end{align*}
Find the solution $u(x,t)$ by separation of variables.
\end{problembox}

\textbf{Assumptions:} Small amplitude ($|u_x| \ll 1$), constant tension $T$ and density $\mu$ ($c = \sqrt{T/\mu}$), ideal flexibility, no external forces or damping.

\textbf{Separation of Variables}
Ansatz $u(x,t) = X(x)T(t)$ yields the coupled ODEs via separation constant $-\lambda$:
\begin{equation}
    \frac{X''}{X} = \frac{1}{c^2}\frac{T''}{T} = -\lambda
\end{equation}

\textbf{Spatial Eigenvalue Problem}
Solving $X'' + \lambda X = 0$ subject to $X(0)=X(L)=0$ imposes $\lambda > 0$.
\begin{itemize}
    \item \textbf{Eigenvalues:} $\lambda_n = \left(\frac{n\pi}{L}\right)^2, \quad n \in \mathbb{N}$
    \item \textbf{Eigenfunctions:} $X_n(x) = \sin\left(\frac{n\pi x}{L}\right)$
\end{itemize}

\textbf{Temporal Evolution}
Solving $T_n'' + \omega_n^2 T_n = 0$ with angular frequency $\omega_n = c\sqrt{\lambda_n} = \frac{n\pi c}{L}$:
\begin{equation}
    T_n(t) = A_n \cos(\omega_n t) + B_n \sin(\omega_n t)
\end{equation}

\begin{summarybox}[title={General Solution}]
By linearity, the general solution is the superposition of normal modes:
$$ u(x,t) = \sum_{n=1}^{\infty} \sin\left(\frac{n\pi x}{L}\right) \left[ A_n \cos\left(\frac{n\pi c t}{L}\right) + B_n \sin\left(\frac{n\pi c t}{L}\right) \right] $$
The Fourier coefficients are determined by orthogonality of sines:
\begin{align*}
    A_n &= \frac{2}{L} \int_0^L f(x) \sin\left(\frac{n\pi x}{L}\right) \, dx \\
    B_n &= \frac{2}{n\pi c} \int_0^L g(x) \sin\left(\frac{n\pi x}{L}\right) \, dx
\end{align*}
\end{summarybox}


\newpage
\subsection{\texorpdfstring{$u_{tt} - c^2 u_{xx} = A\sin\!\left(\frac{3\pi x}{L}\right)\cos(\Omega t)$ on $(0,L)$; $u(0,t)=u(L,t)=0$; zero ICs; resonant ($\Omega=\omega_3$) and non-resonant cases}{Forced wave equation: bounded beats when non-resonant; secular growth when resonant}}

\begin{problembox}
$$u_{tt} - c^2 u_{xx} = A \sin\left(\frac{3\pi x}{L}\right)\cos(\Omega t)$$
BCs: $u(0, t) = u(L, t) = 0$. ICs: $u(x, 0) = 0, u_t(x, 0) = 0$.
\end{problembox}


We use the method of eigenfunction expansion. The eigenfunctions of the homogeneous operator with zero BCs are $\sin\left(\frac{n\pi x}{L}\right)$.
Ansatz:
$$ u(x,t) = \sum_{n=1}^{\infty} T_n(t) \sin\left(\frac{n\pi x}{L}\right) $$
Substituting into the PDE:
$$ \sum_{n=1}^{\infty} \left[ T_n''(t) + \left(\frac{n\pi c}{L}\right)^2 T_n(t) \right] \sin\left(\frac{n\pi x}{L}\right) = A \sin\left(\frac{3\pi x}{L}\right)\cos(\Omega t) $$
By orthogonality, with natural frequency $\omega_n = \tfrac{n\pi c}{L}$, we equate coefficients for each $n$.
For $n \neq 3$:
$$ T_n'' + \omega_n^2 T_n = 0, \quad T_n(0)=0, T_n'(0)=0 \implies T_n(t) = 0 $$
For $n = 3$: $T_3'' + \omega_3^2 T_3 = A\cos(\Omega t)$, with $T_3(0)=0$, $T_3'(0)=0$.

\textbf{Case 1: Non-Resonant ($\Omega \neq \omega_3$)}
The general solution is $T_3(t) = c_1 \cos(\omega_3 t) + c_2 \sin(\omega_3 t) + \frac{A}{\omega_3^2 - \Omega^2}\cos(\Omega t)$.
$T_3'(0) = \omega_3 c_2 = 0 \implies c_2 = 0$.
$T_3(0) = c_1 + \frac{A}{\omega_3^2 - \Omega^2} = 0 \implies c_1 = - \frac{A}{\omega_3^2 - \Omega^2}$.
$$ T_3(t) = \frac{A}{\omega_3^2 - \Omega^2} (\cos(\Omega t) - \cos(\omega_3 t)) $$

\textbf{Case 2: Resonant ($\Omega = \omega_3$)}
The particular solution form is $t(a \sin \Omega t + b \cos \Omega t)$.
Solving $T_3'' + \Omega^2 T_3 = A \cos(\Omega t)$ yields particular solution $T_p = \frac{A}{2\Omega} t \sin(\Omega t)$.
Homogeneous part satisfies initial conditions (starts at 0 with 0 slope).
$$ T_3(t) = \frac{A}{2\omega_3} t \sin(\omega_3 t) $$

\begin{summarybox}
Solution: $u(x,t) = T_3(t) \sin\left(\frac{3\pi x}{L}\right)$.
\begin{itemize}
    \item If $\Omega \neq \frac{3\pi c}{L}$: $T_3(t) = \frac{A}{\omega_3^2 - \Omega^2} (\cos(\Omega t) - \cos(\omega_3 t))$
    \item If $\Omega = \frac{3\pi c}{L}$: $T_3(t) = \frac{A L}{6\pi c} t \sin\left(\frac{3\pi c t}{L}\right)$
\end{itemize}
\end{summarybox}

\newpage
\subsection{$\frac{\partial^2 u}{\partial t^2}-\frac{\partial^2 u}{\partial x^2}=0,\quad u(x,0)=f(x),\quad u_t(x,0)=g(x)$}
\begin{problembox}
$$\frac{\partial^2 u}{\partial t^2}-\frac{\partial^2 u}{\partial x^2}=0, \quad x\in[0,\pi],\; t>0$$
BCs: $u(0,t)=0,\quad u(\pi,t)=0$.

(a) Find the solution with ICs: $u(x,0)=f(x),\quad u_t(x,0)=g(x)$.

(b) Replace the second IC with $u(x,t_0) = h(x)$ for some $t_0 > 0$. Under what condition on $t_0$ does a unique solution exist?
\end{problembox}

\textbf{(a) Solution by separation of variables.}

Assume $u(x,t)=X(x)T(t)$. Substituting into the PDE and separating:
$$\frac{T''(t)}{T(t)}=\frac{X''(x)}{X(x)}=-\lambda$$
This yields the spatial ODE $X''(x)+\lambda X(x)=0$ and the temporal ODE $T''(t)+\lambda T(t)=0$.

Nontrivial solutions of the spatial problem exist only for
\[
\lambda=n^2, \qquad n\in\mathbb{N},
\]
with eigenfunctions
\[
X_n(x)=\sin(nx).
\]

For each $n$,
\[
T_n(t)=A_n\cos(nt)+B_n\sin(nt).
\]

By superposition,
\[
u(x,t)=\sum_{n=1}^{\infty}
\left[A_n\cos(nt)+B_n\sin(nt)\right]\sin(nx).
\]

Imposing $u(x,0)=f(x)$ gives
\[
f(x)=\sum_{n=1}^{\infty}A_n\sin(nx),
\qquad
A_n=\frac{2}{\pi}\int_0^{\pi}f(x)\sin(nx)\,dx.
\]

Imposing $u_t(x,0)=g(x)$ gives
\[
g(x)=\sum_{n=1}^{\infty}nB_n\sin(nx),
\qquad
B_n=\frac{2}{n\pi}\int_0^{\pi}g(x)\sin(nx)\,dx.
\]

Thus
\[
u(x,t)=\sum_{n=1}^{\infty}
\left[
\left(\frac{2}{\pi}\int_0^{\pi}f(x)\sin(nx)\,dx\right)\cos(nt)
+
\left(\frac{2}{n\pi}\int_0^{\pi}g(x)\sin(nx)\,dx\right)\sin(nt)
\right]\sin(nx).
\]

\textbf{(b) Modified initial condition.}

The general form remains
\[
u(x,t)=\sum_{n=1}^{\infty}
\left[A_n\cos(nt)+B_n\sin(nt)\right]\sin(nx),
\]
\[
A_n=\frac{2}{\pi}\int_0^{\pi}f(x)\sin(nx)\,dx.
\]

Define the sine coefficients of $h(x)$:
\[
H_n=\frac{2}{\pi}\int_0^{\pi}h(x)\sin(nx)\,dx.
\]

The condition $u(x,t_0)=h(x)$ yields
\[
A_n\cos(nt_0)+B_n\sin(nt_0)=H_n,
\]
hence
\[
B_n=\frac{H_n-A_n\cos(nt_0)}{\sin(nt_0)}.
\]

A unique solution exists if and only if
\[
\sin(nt_0)\neq0
\quad \forall n\in\mathbb{N},
\]
equivalently
\[
t_0\neq\frac{m}{n}\pi
\quad\Longleftrightarrow\quad
\frac{t_0}{\pi}\notin\mathbb{Q}.
\]

The solution is therefore
\[
u(x,t)=\sum_{n=1}^{\infty}
\left[
A_n\cos(nt)
+
\frac{H_n-A_n\cos(nt_0)}{\sin(nt_0)}\,\sin(nt)
\right]\sin(nx).
\]

\newpage
% ---------------------------------------------------------------------
% NEW PROBLEM (Section 3)
% ---------------------------------------------------------------------
\subsection{\texorpdfstring{$u_{tt}=c^2 u_{xx}$ on $(-\infty,\infty)$; derive d'Alembert's formula; apply to $u(x,0)=e^{-x^2}$, $u_t(x,0)=0$}{D'Alembert's formula via characteristic coordinates; splitting into two traveling Gaussian pulses}}

\begin{problembox}
Consider the wave equation on the full real line:
$$u_{tt} = c^2 u_{xx}, \quad -\infty < x < \infty, \; t > 0$$

(a) Derive d'Alembert's formula for the Cauchy problem with initial data $u(x,0) = \phi(x)$ and $u_t(x,0) = \psi(x)$.

(b) Find the explicit solution when $\phi(x) = e^{-x^2}$ and $\psi(x) = 0$. Describe the physical behaviour for large $t$.
\end{problembox}

\textbf{(a) Derivation of d'Alembert's Formula}

\textbf{Step 1: Change to characteristic coordinates.}
Introduce $\xi = x + ct$ and $\eta = x - ct$. Then:
\[
    u_t = c\,(u_\xi - u_\eta), \quad
    u_{tt} = c^2(u_{\xi\xi} - 2u_{\xi\eta} + u_{\eta\eta}),
\]
\[
    u_x = u_\xi + u_\eta, \quad
    u_{xx} = u_{\xi\xi} + 2u_{\xi\eta} + u_{\eta\eta}.
\]
Substituting into the PDE:
\[
    c^2(u_{\xi\xi} - 2u_{\xi\eta} + u_{\eta\eta}) = c^2(u_{\xi\xi} + 2u_{\xi\eta} + u_{\eta\eta})
    \implies u_{\xi\eta} = 0.
\]

\textbf{Step 2: Integrate the canonical form.}
$u_{\xi\eta} = 0$ integrates immediately to:
\[
    u(\xi,\eta) = F(\xi) + G(\eta) \implies u(x,t) = F(x+ct) + G(x-ct),
\]
for arbitrary $C^1$ functions $F$ and $G$.

\textbf{Step 3: Apply initial conditions.}
At $t=0$:
\begin{align*}
    u(x,0) &= F(x) + G(x) = \phi(x), \tag{i}\\
    u_t(x,0) &= c\big[F'(x) - G'(x)\big] = \psi(x). \tag{ii}
\end{align*}
Integrating (ii):
\[
    F(x) - G(x) = \frac{1}{c}\int_0^x \psi(s)\,ds + K. \tag{iii}
\]
Adding (i) and (iii) and subtracting (i) from (iii):
\[
    F(x) = \frac{1}{2}\phi(x) + \frac{1}{2c}\int_0^x\psi(s)\,ds + \frac{K}{2}, \qquad
    G(x) = \frac{1}{2}\phi(x) - \frac{1}{2c}\int_0^x\psi(s)\,ds - \frac{K}{2}.
\]
Substituting back and noting the constants $\pm K/2$ cancel:

\begin{summarybox}[title={D'Alembert's Formula}]
$$u(x,t) = \frac{\phi(x+ct) + \phi(x-ct)}{2} + \frac{1}{2c}\int_{x-ct}^{x+ct}\psi(s)\,ds$$
\end{summarybox}

\textbf{(b) Gaussian Initial Displacement}

With $\phi(x) = e^{-x^2}$ and $\psi(x) = 0$, d'Alembert's formula gives directly:

\begin{summarybox}
$$u(x,t) = \frac{1}{2}e^{-(x+ct)^2} + \frac{1}{2}e^{-(x-ct)^2}$$
\end{summarybox}

\textbf{Physical interpretation.}
The initial Gaussian pulse at $x=0$ splits at $t=0^+$ into two identical pulses of half amplitude:
\begin{itemize}
    \item A right-travelling pulse centred at $x = ct$: $\frac{1}{2}e^{-(x-ct)^2}$.
    \item A left-travelling pulse centred at $x = -ct$: $\frac{1}{2}e^{-(x+ct)^2}$.
\end{itemize}
For large $t$, the two pulses are well separated (centres at $\pm ct$), and the region between them ($|x| \ll ct$) is quiescent ($u \approx 0$). The shape of each pulse is preserved identically since the wave equation is non-dispersive.

\newpage

% ---------------------------------------------------------------------
% PROBLEM 5
% ---------------------------------------------------------------------
\section{Laplace and Poisson Equations}

\begin{notebox}[title={General Procedure: Laplace \& Poisson Equations}]
For $\nabla^2 u = 0$ (Laplace) or $\nabla^2 u = S$ (Poisson) in a given domain, the procedure depends on the geometry.

\textbf{General Steps:}
\begin{enumerate}
    \item \textbf{Choose coordinates} matching the domain: Cartesian for rectangles, polar $(r,\theta)$ for disks, annuli, and sectors.
    \item \textbf{For Poisson equations:} decompose $u = u_p + u_h$, where $\nabla^2 u_p = S$ (find a particular solution, often radial) and $\nabla^2 u_h = 0$ with modified BCs $u_h|_{\text{boundary}} = u|_{\text{boundary}} - u_p|_{\text{boundary}}$.
    \item \textbf{For non-homogeneous BCs on straight edges:} split $u = v + w$, where $v$ absorbs the non-homogeneous data (often a linear/polynomial function satisfying Laplace) and $w$ has homogeneous BCs on those edges.
    \item \textbf{Separate variables.} In polar: $u=R(r)\Theta(\theta)$. The angular part produces eigenfunctions ($\sin, \cos$); the radial part satisfies an Euler--Cauchy equation with solutions $r^{\pm n}$ (or $\ln r$ for $n=0$).
    \item \textbf{Enforce regularity/boundedness:} exclude $r^{-n}$ at the origin, $\ln r$ at the origin, etc.
    \item \textbf{Apply remaining BCs.} Match the boundary data by expanding in eigenfunctions and equating Fourier coefficients.
\end{enumerate}
\textbf{Key domain types:}
\begin{itemize}
    \item \textbf{Disk / sector:} $R_n = r^n$ (bounded at origin); expand the BC at $r=R$ in $\sin/\cos$.
    \item \textbf{Annulus:} $R_n = A_n r^n + B_n r^{-n}$; two radial BCs determine both constants.
    \item \textbf{Rectangle:} Eigenfunctions in one variable ($\sin$), $\sinh/\cosh$ in the other; use the three homogeneous BCs to constrain the form, then match the non-homogeneous BC.
\end{itemize}
\end{notebox}

\subsection{\texorpdfstring{$\nabla^2 u = 0$ in a circular sector of radius $R$ and angle $\alpha$}{Laplace in sector}}

\begin{problembox}
Solve $\nabla^2 u=0$ in $0<r<R$, $0<\theta<\alpha$ with
\[
u(r,0)=u_0,\quad
u(r,\alpha)=0,\quad
u(R,\theta)=0,
\]
where $u_0\neq0$ is constant. Require boundedness as $r\to0$.
BCs: Dirichlet on $\theta=0,\alpha$ and $r=R$.
\end{problembox}

In polar coordinates,
\[
\nabla^2 u=u_{rr}+\frac1r u_r+\frac1{r^2}u_{\theta\theta}.
\]

Since $u(r,0)\neq0$, write $u=v(\theta)+w(r,\theta)$ with $v''(\theta)=0$, hence
$v(\theta)=A\theta+B$. From $v(0)=u_0$, $v(\alpha)=0$ we obtain
$B=u_0$, $A=-u_0/\alpha$, so
\[
v(\theta)=u_0\!\left(1-\frac{\theta}{\alpha}\right).
\]

Then $w$ satisfies $\nabla^2 w=0$ with
$w(r,0)=w(r,\alpha)=0$ and
$w(R,\theta)=-u_0(1-\theta/\alpha)$.

Seek $w=\mathcal R(r)\Theta(\theta)$. Substitution gives
$r^2\mathcal R''+r\mathcal R'-\lambda\mathcal R=0$ and
$\Theta''+\lambda\Theta=0$.

From $\Theta(0)=\Theta(\alpha)=0$ we need
$\lambda_n=(n\pi/\alpha)^2$ and
$\Theta_n(\theta)=\sin(n\pi\theta/\alpha)$.

For the radial part, try $\mathcal R=r^p$:
$p(p-1)+p-\lambda_n=0 \Rightarrow p^2-\lambda_n=0$,
so $p=\pm n\pi/\alpha$.
Boundedness at $r\to0$ excludes the negative power, hence
$\mathcal R_n(r)=r^{n\pi/\alpha}$.

Thus
\[
w(r,\theta)=\sum_{n=1}^\infty
A_n r^{n\pi/\alpha}
\sin\!\left(\frac{n\pi\theta}{\alpha}\right).
\]

Imposing $r=R$ gives
\[
\sum_{n=1}^\infty
A_n R^{n\pi/\alpha}
\sin\!\left(\frac{n\pi\theta}{\alpha}\right)
=
-\,u_0\!\left(1-\frac{\theta}{\alpha}\right).
\]

Expand the right-hand side in a sine series on $[0,\alpha]$:
\[
C_n=\frac{2}{\alpha}\int_0^\alpha
\!\left[-u_0\!\left(1-\frac{\theta}{\alpha}\right)\right]
\sin\!\left(\frac{n\pi\theta}{\alpha}\right)d\theta.
\]

Compute
\[
I_n=\int_0^\alpha
\left(1-\frac{\theta}{\alpha}\right)
\sin\!\left(\frac{n\pi\theta}{\alpha}\right)d\theta.
\]
Let $x=\frac{\pi\theta}{\alpha}$ so
$I_n=\frac{\alpha}{\pi}\int_0^\pi(1-x/\pi)\sin(nx)\,dx$.
Using
$\int_0^\pi\sin(nx)\,dx=\frac{2}{n}(1-(-1)^n)$ and
$\int_0^\pi x\sin(nx)\,dx=\frac{\pi(-1)^{n+1}}{n}$,
one finds $I_n=\alpha/(n\pi)$.
Hence $C_n=-2u_0/(n\pi)$.

Since $C_n=A_n R^{n\pi/\alpha}$,
$A_n=-\frac{2u_0}{n\pi}R^{-n\pi/\alpha}$.

Therefore
\[
w(r,\theta)=
-\frac{2u_0}{\pi}
\sum_{n=1}^\infty
\frac1n
\left(\frac{r}{R}\right)^{n\pi/\alpha}
\sin\!\left(\frac{n\pi\theta}{\alpha}\right).
\]

\begin{summarybox}
\[
u(r,\theta)=
u_0\!\left(1-\frac{\theta}{\alpha}\right)
-
\frac{2u_0}{\pi}
\sum_{n=1}^\infty
\frac1n
\left(\frac{r}{R}\right)^{n\pi/\alpha}
\sin\!\left(\frac{n\pi\theta}{\alpha}\right).
\]
\end{summarybox}
\subsection{\texorpdfstring{$\nabla^2 u = 0$ in semi-annulus $1<r<2$, $0<\theta<\pi$; $u(r,0)=u(r,\pi)=0$, $u(1,\theta)=0$, $u(2,\theta)=5\sin 2\theta-\sin 4\theta$}{Laplace in semi-annulus: zero on straight edges and inner arc, sine data on outer arc}}

\begin{problembox}
Solve $\nabla^2 u = 0$ in the semi-annulus $1 < r < 2, 0 < \theta < \pi$.
BCs: $u(r, 0) = u(r, \pi) = 0$, $u(1, \theta) = 0$,
$u(2, \theta) = 5\sin(2\theta) - \sin(4\theta)$.
\end{problembox}

In polar coordinates, $\nabla^2 u = u_{rr} + \frac{1}{r}u_r + \frac{1}{r^2}u_{\theta\theta} = 0$.
Separation of variables $u(r,\theta) = R(r)\Theta(\theta)$ with homogeneous angular BCs $\Theta(0)=\Theta(\pi)=0$ leads to:
$$ \Theta_n(\theta) = \sin(n\theta), \quad n = 1, 2, 3, \dots $$
The radial equation is Euler-Cauchy: $r^2 R'' + r R' - n^2 R = 0$.
Solution: $R_n(r) = A_n r^n + B_n r^{-n}$.
General solution:
$$ u(r, \theta) = \sum_{n=1}^{\infty} (A_n r^n + B_n r^{-n}) \sin(n\theta) $$

\textbf{Apply Radial BCs:}

\textbf{1. Inner BC:} $u(1, \theta) = 0$:
$$ \sum_{n=1}^{\infty} (A_n + B_n) \sin(n\theta) = 0 \implies B_n = -A_n $$
So $R_n(r) = A_n(r^n - r^{-n})$, and:
$$ u(r, \theta) = \sum_{n=1}^{\infty} A_n (r^n - r^{-n}) \sin(n\theta) $$

\textbf{2. Outer BC:} $u(2, \theta) = 5\sin(2\theta) - \sin(4\theta)$:
$$ \sum_{n=1}^{\infty} A_n (2^n - 2^{-n}) \sin(n\theta) = 5\sin(2\theta) - \sin(4\theta) $$
Matching coefficients for each mode:

\textbf{For $n=2$:} $A_2(2^2 - 2^{-2}) = A_2 \cdot \tfrac{15}{4} = 5 \implies A_2 = \tfrac{4}{3}$.

\textbf{For $n=4$:} $A_4(2^4 - 2^{-4}) = A_4 \cdot \tfrac{255}{16} = -1 \implies A_4 = -\tfrac{16}{255}$.

All other $A_n = 0$.

\begin{summarybox}
$$ u(r, \theta) = \frac{4}{3}\left(r^2 - \frac{1}{r^2}\right)\sin(2\theta) - \frac{16}{255}\left(r^4 - \frac{1}{r^4}\right)\sin(4\theta) $$
\end{summarybox}

\newpage
\subsection{\texorpdfstring{$\nabla^2 u = 0$ in annulus $1<r<2$; (A) $u(1)=A\sin\theta$, $u(2)=B\sin 2\theta$; (B) $u(1)=0$, $u_r(2)=2\cos\theta$}{Laplace in full annulus: (A) two-mode Dirichlet BCs; (B) mixed Dirichlet--Neumann BCs}}

\begin{problembox}
Solve the Laplace equation $\nabla^2 u = 0$ in the annular domain $D = \{ (r, \theta) \mid 1 < r < 2, \ 0 \le \theta < 2\pi \}$.

(A) Dirichlet BCs: $u(1, \theta) = A\sin\theta$, \quad $u(2, \theta) = B\sin(2\theta)$.

(B) Mixed BCs: $u(1,\theta) = 0$, $\quad \left. \dfrac{\partial u}{\partial r} \right|_{r=2} = 2\cos\theta$
\end{problembox}
In polar coordinates, $\nabla^2 u = u_{rr} + \frac{1}{r}u_r + \frac{1}{r^2}u_{\theta\theta} = 0$.
Separation of variables $u(r,\theta) = R(r)\Theta(\theta)$ leads to $\Theta_n(\theta) = A_n\cos(n\theta) + B_n\sin(n\theta)$, $n = 0,1,2,\dots$

The radial equation is Euler-Cauchy: $r^2 R'' + r R' - n^2 R = 0$.
For $n \neq 0$, the solution is $R_n(r) = c_n r^n + d_n r^{-n}$.
For $n = 0$, the solution is $R_0(r) = c_0 + d_0 \ln r$.

The general solution in an annulus is:
$$ u(r, \theta) = c_0 + d_0 \ln r + \sum_{n=1}^{\infty} \left( c_n r^n + d_n r^{-n} \right) (A_n \cos(n\theta) + B_n \sin(n\theta)) $$

\textbf{(A) Dirichlet Boundary Conditions}

Given BCs: $u(1, \theta) = A\sin\theta$ and $u(2, \theta) = B\sin(2\theta)$.

Since the boundary data contains only $\sin\theta$ and $\sin(2\theta)$ modes, we take $u(r, \theta) = R_1(r)\sin\theta + R_2(r)\sin(2\theta)$.

\textbf{Step 1: The $n=1$ mode.}
Let $R_1(r) = a_1 r + b_1 r^{-1}$. The boundary conditions require:
\begin{align*}
    R_1(1) &= a_1 + b_1 = A \\
    R_1(2) &= 2a_1 + \frac{b_1}{2} = 0
\end{align*}
From the second equation: $b_1 = -4a_1$. Substituting into the first:
$$ a_1 - 4a_1 = A \implies a_1 = -\frac{A}{3}, \quad b_1 = \frac{4A}{3} $$

\textbf{Step 2: The $n=2$ mode.}
Let $R_2(r) = a_2 r^2 + b_2 r^{-2}$. The boundary conditions require:
\begin{align*}
    R_2(1) &= a_2 + b_2 = 0 \implies b_2 = -a_2 \\
    R_2(2) &= 4a_2 + \frac{b_2}{4} = B
\end{align*}
Substituting $b_2 = -a_2$ into the second equation:
$$ 4a_2 - \frac{a_2}{4} = \frac{15}{4}a_2 = B \implies a_2 = \frac{4B}{15}, \quad b_2 = -\frac{4B}{15} $$

\begin{summarybox}[title={Part A Answer}]
$$ u(r, \theta) = \frac{A}{3} \left( \frac{4}{r} - r \right) \sin\theta + \frac{4B}{15} \left( r^2 - \frac{1}{r^2} \right) \sin(2\theta) $$
\end{summarybox}

\textbf{(B) Mixed Boundary Conditions}

Given BCs: $u(1, \theta) = 0$ and $\left.\frac{\partial u}{\partial r}\right|_{r=2} = 2\cos\theta$.

Since the Neumann data at $r=2$ contains only a $\cos\theta$ mode, the solution involves only the $n=1$ cosine mode:
$$ u(r, \theta) = R(r) \cos\theta, \quad \text{with } R(r) = c_1 r + c_2 r^{-1} $$

\textbf{Step 1: Apply inner Dirichlet BC ($r=1$).}
$$ R(1) = c_1 + c_2 = 0 \implies c_2 = -c_1 $$
Thus $R(r) = c_1(r - r^{-1})$ and $R'(r) = c_1(1 + r^{-2})$.

\textbf{Step 2: Apply outer Neumann BC ($r=2$).}
$$ R'(2) = c_1 \left( 1 + \frac{1}{4} \right) = \frac{5c_1}{4} = 2 \implies c_1 = \frac{8}{5} $$

\begin{summarybox}[title={Part B Answer}]
$$ u(r, \theta) = \frac{8}{5} \left( r - \frac{1}{r} \right) \cos\theta $$
\end{summarybox}
\subsection{\texorpdfstring{$\nabla^2 u = 0$ in sector $0\le\theta\le\pi$, $r\le R$; $u(r,0)=u(r,\pi)=0$, $u(R,\theta)=2\sin\theta+4\sin 3\theta$}{Laplace in half-disk: Dirichlet on straight edges, two-mode sine data on arc}}

\begin{problembox}
Find a harmonic function $u(r,\theta)$ in the sector $0 \le \theta \le \pi$ with boundary conditions:
\begin{itemize}
    \item $u(R, \theta) = 2\sin\theta + 4\sin 3\theta$ on the circular arc
    \item $u(r, 0) = u(r, \pi) = 0$ on the straight edges
\end{itemize}
\end{problembox}

\textbf{Solution:}

Separation of variables $u(r,\theta) = R(r)\Theta(\theta)$ with $\Theta(0)=\Theta(\pi)=0$ yields eigenfunctions:
$$ \Theta_n(\theta) = \sin(n\theta), \quad n=1,2,\dots $$
The corresponding radial solutions bounded at origin are $R_n(r) = r^n$.
Thus, the general solution is:
$$ u(r,\theta) = \sum_{n=1}^{\infty} A_n r^n \sin(n\theta) $$

Boundary condition at $r=R$:
$$ u(R,\theta) = \sum_{n=1}^{\infty} A_n R^n \sin(n\theta) = 2\sin\theta + 4\sin 3\theta $$

Comparing coefficients of $\sin(n\theta)$:
\begin{itemize}
    \item For $n=1$: $A_1 R^1 = 2 \implies A_1 = \frac{2}{R}$
    \item For $n=3$: $A_3 R^3 = 4 \implies A_3 = \frac{4}{R^3}$
    \item For all other $n$: $A_n = 0$
\end{itemize}

Substituting these back into the general solution:

\begin{summarybox}
$$ u(r, \theta) = \frac{2r}{R}\sin\theta + \frac{4r^3}{R^3}\sin 3\theta $$
\end{summarybox}

\newpage

% ---------------------------------------------------------------------
% PROBLEM 13
% ---------------------------------------------------------------------
\subsection{\texorpdfstring{$\nabla^2 u = r$ inside unit circle $r<1$; $u(1,\theta)=0$}{Poisson with radial source: radially symmetric ODE, regularity at origin forces $C_1=0$}}

\begin{problembox}
Solve the Poisson equation $\nabla^2 u = r$ inside the unit circle $r < 1$ with boundary condition $u(1, \theta) = 0$.
\end{problembox}

\textbf{Solution:}

Since the source term $r$ and the boundary condition (zero on circle) are rotationally invariant (independent of $\theta$), we seek a solution $u(r)$ depending only on $r$.
The Laplacian in polar coordinates for radial functions is:
$$ \frac{1}{r} \frac{d}{dr}\left(r \frac{du}{dr}\right) = r $$

Integrate w.r.t $r$:
$$ \frac{d}{dr}\left(r \frac{du}{dr}\right) = r^2 $$
$$ r \frac{du}{dr} = \frac{r^3}{3} + C_1 $$
$$ \frac{du}{dr} = \frac{r^2}{3} + \frac{C_1}{r} $$

Integrate again:
$$ u(r) = \frac{r^3}{9} + C_1 \ln r + C_2 $$

\textbf{Boundary Conditions:}
1. \textbf{Regularity at origin:} $u(0)$ must be finite. Since $\ln r \to -\infty$, we must set $C_1 = 0$.
2. \textbf{Boundary Value:} $u(1) = \tfrac{1}{9} + C_2 = 0 \implies C_2 = -\tfrac{1}{9}$.

\begin{summarybox}
$$ u(r) = \frac{1}{9}(r^3 - 1) $$
\end{summarybox}
\newpage
\subsection{\texorpdfstring{$\nabla^2 u = 16r^2$ in annulus $1<r<2$; $u|_{r=1}=3\sin\theta$, $u|_{r=2}=12\cos\theta$}{Poisson in annulus: particular radial solution $u_p=r^4$ plus homogeneous Laplace for boundary residuals}}

\begin{problembox}
Solve the Poisson equation $\frac{\partial^2 u}{\partial x^2} + \frac{\partial^2 u}{\partial y^2} = 16(x^2 + y^2)$ in the annular domain $D = \{ (r, \theta) \mid 1 < r < 2, \ 0 \le \theta < 2\pi \}$, subject to the boundary conditions:
$$ u|_{r=2} = 6x, \quad u|_{r=1} = 3y $$
Hint: The general solution of the Poisson equation can be written as the sum of a particular solution and the general solution of the Laplace equation.
\end{problembox}

The problem is defined on an annulus with a radially symmetric source term. We switch to polar coordinates where $x = r\cos\theta$, $y = r\sin\theta$, and $x^2+y^2 = r^2$.
The equation becomes:
$$ \nabla^2 u = \frac{1}{r} \frac{\partial}{\partial r} \left( r \frac{\partial u}{\partial r} \right) + \frac{1}{r^2} \frac{\partial^2 u}{\partial \theta^2} = 16r^2 $$

We decompose $u(r,\theta) = u_p(r) + u_h(r,\theta)$ where $\nabla^2 u_p = 16r^2$ and $\nabla^2 u_h = 0$.

\textbf{Step 1: The Particular Solution ($u_p$)}

Since the source term $S(r) = 16r^2$ is independent of $\theta$, we assume a radial particular solution $u_p = u_p(r)$. The Laplacian reduces to the radial part:
$$ \frac{1}{r} \frac{d}{dr} \left( r \frac{du_p}{dr} \right) = 16r^2 $$
Multiply by $r$ and integrate:
$$ \frac{d}{dr} \left( r \frac{du_p}{dr} \right) = 16r^3 \implies r \frac{du_p}{dr} = 4r^4 + C_1 $$
Divide by $r$ and integrate again:
$$ \frac{du_p}{dr} = 4r^3 + \frac{C_1}{r} \implies u_p(r) = r^4 + C_1 \ln r + C_2 $$
Since $C_1 \ln r + C_2$ satisfies the homogeneous Laplace equation, we set $C_1=C_2=0$, giving $u_p(r) = r^4$.

\textbf{Step 2: The Homogeneous Solution ($u_h$)}

The homogeneous function $u_h$ must satisfy $\nabla^2 u_h = 0$ with modified boundary conditions.
The total solution is $u = r^4 + u_h$.
We convert the given BCs to polar coordinates and subtract $u_p$:
\begin{align*}
    \text{At } r=1: \quad & u(1, \theta) = 3(1)\sin\theta = 3\sin\theta \\
    & u_h(1, \theta) = u(1, \theta) - u_p(1) = 3\sin\theta - 1^4 = 3\sin\theta - 1 \\
    \text{At } r=2: \quad & u(2, \theta) = 6(2)\cos\theta = 12\cos\theta \\
    & u_h(2, \theta) = u(2, \theta) - u_p(2) = 12\cos\theta - 2^4 = 12\cos\theta - 16
\end{align*}

The general solution for the Laplacian in an annulus is:
$$ u_h(r, \theta) = A_0 + B_0 \ln r + \sum_{n=1}^{\infty} (A_n r^n + B_n r^{-n})\cos(n\theta) + (C_n r^n + D_n r^{-n})\sin(n\theta) $$
By inspection of the boundary terms ($1, \sin\theta, \cos\theta$), we only need the isotropic ($n=0$) and $n=1$ modes.

\textbf{2a. Isotropic Mode ($n=0$)}
Let $H_0(r) = A_0 + B_0 \ln r$. Matching the constant terms on the boundaries:
\begin{align*}
    H_0(1) &= A_0 = -1 \\
    H_0(2) &= A_0 + B_0 \ln 2 = -16
\end{align*}
Substituting $A_0 = -1$ into the second equation:
$$ -1 + B_0 \ln 2 = -16 \implies B_0 \ln 2 = -15 \implies B_0 = -\frac{15}{\ln 2} $$

\textbf{2b. Cosine Mode ($n=1$)}
Let $H_{\cos}(r) = \left( A_1 r + B_1 r^{-1} \right) \cos\theta$. Matching the $\cos\theta$ terms:
\begin{align*}
    r=1: \quad & A_1 + B_1 = 0 \implies B_1 = -A_1 \\
    r=2: \quad & 2A_1 + \frac{B_1}{2} = 12
\end{align*}
Substituting $B_1 = -A_1$:
$$ 2A_1 - \frac{A_1}{2} = \frac{3}{2}A_1 = 12 \implies A_1 = 8, \quad B_1 = -8 $$

\textbf{2c. Sine Mode ($n=1$)}
Let $H_{\sin}(r) = \left( C_1 r + D_1 r^{-1} \right) \sin\theta$. Matching the $\sin\theta$ terms:
\begin{align*}
    r=1: \quad & C_1 + D_1 = 3 \\
    r=2: \quad & 2C_1 + \frac{D_1}{2} = 0 \implies D_1 = -4C_1
\end{align*}
Substituting into the first equation:
$$ C_1 - 4C_1 = -3C_1 = 3 \implies C_1 = -1, \quad D_1 = 4 $$

\begin{summarybox}[title={Final Solution}]
Combining the particular solution and the calculated homogeneous modes:
$$ u(r, \theta) = r^4 - 1 - \frac{15}{\ln 2} \ln r + 8\left(r - \frac{1}{r}\right) \cos\theta - \left(r - \frac{4}{r}\right) \sin\theta $$
\end{summarybox}

% ---------------------------------------------------------------------
% NEW PROBLEM (Section 4)
% ---------------------------------------------------------------------
\newpage
\subsection{\texorpdfstring{$\nabla^2 u=0$ on rectangle $0<x<\pi,\,0<y<\pi$; zero on three sides, $u(x,\pi)=\sin x+3\sin 2x$}{Laplace on rectangle: eigenfunction expansion in $x$; sinh solution in $y$ satisfying zero bottom BC}}

\begin{problembox}
Solve $\nabla^2 u = 0$ on the rectangle $\{0 < x < \pi,\; 0 < y < \pi\}$ subject to:
$$u(0,y) = 0, \quad u(\pi,y) = 0, \quad u(x,0) = 0, \quad u(x,\pi) = \sin x + 3\sin 2x.$$
\end{problembox}

\textbf{Step 1: Separation of variables.}
Writing $u(x,y) = X(x)Y(y)$ and substituting into $u_{xx} + u_{yy} = 0$:
\[
    \frac{X''}{X} = -\frac{Y''}{Y} = -\lambda.
\]
The homogeneous BCs in $x$ give the Sturm-Liouville problem:
\[
    X'' + \lambda X = 0, \quad X(0) = X(\pi) = 0
    \implies \lambda_n = n^2, \quad X_n(x) = \sin(nx), \quad n = 1,2,3,\dots
\]

\textbf{Step 2: Solve the $y$-equation.}
For each $n$: $Y'' - n^2 Y = 0$, with general solution $Y_n(y) = A_n \sinh(ny) + B_n \cosh(ny)$.
The BC $u(x,0)=0$ requires $Y_n(0) = 0$:
\[
    B_n \cosh(0) = B_n = 0.
\]
Hence $Y_n(y) = A_n \sinh(ny)$, and the general solution is:
\[
    u(x,y) = \sum_{n=1}^{\infty} C_n \sinh(ny)\sin(nx).
\]

\textbf{Step 3: Apply the non-homogeneous BC at $y = \pi$.}
\[
    u(x,\pi) = \sum_{n=1}^{\infty} C_n \sinh(n\pi)\sin(nx) = \sin x + 3\sin 2x.
\]
Comparing coefficients of $\sin(nx)$:
\begin{itemize}
    \item $n=1$: $C_1 \sinh(\pi) = 1 \implies C_1 = \dfrac{1}{\sinh\pi}$
    \item $n=2$: $C_2 \sinh(2\pi) = 3 \implies C_2 = \dfrac{3}{\sinh 2\pi}$
    \item $n \ge 3$: $C_n = 0$
\end{itemize}

\begin{summarybox}
$$u(x,y) = \frac{\sinh y}{\sinh\pi}\sin x + \frac{3\sinh 2y}{\sinh 2\pi}\sin 2x$$
\end{summarybox}

\textbf{Verification:}
\begin{itemize}
    \item $u(x,0) = 0$ (since $\sinh 0 = 0$). \checkmark
    \item $u(0,y) = u(\pi,y) = 0$ (since $\sin 0 = \sin(n\pi) = 0$). \checkmark
    \item $u(x,\pi) = \frac{\sinh\pi}{\sinh\pi}\sin x + \frac{3\sinh2\pi}{\sinh2\pi}\sin 2x = \sin x + 3\sin 2x$. \checkmark
    \item $\nabla^2 u = 0$: each term $\frac{\sinh(ny)}{\sinh(n\pi)}\sin(nx)$ satisfies Laplace since $(-n^2)\sin(nx)\cdot\frac{\sinh(ny)}{\sinh n\pi} + n^2\sin(nx)\cdot\frac{\sinh(ny)}{\sinh n\pi} = 0$. \checkmark
\end{itemize}

\newpage
% ---------------------------------------------------------------------
% PROBLEM 6
% ---------------------------------------------------------------------
\section{Calculus of Variations}

\begin{notebox}[title={General Procedure: Calculus of Variations}]
Given a functional $I[y]=\int_a^b L(x,y,y')\,dx$ to extremize, subject to boundary conditions $y(a)=\alpha$, $y(b)=\beta$ and possibly integral constraints:
\begin{enumerate}
    \item \textbf{Write the Euler--Lagrange (E-L) equation:}
    $\displaystyle\frac{\partial L}{\partial y} - \frac{d}{dx}\frac{\partial L}{\partial y'} = 0.$
    \item \textbf{If $L$ is independent of $x$}, use the \textbf{Beltrami identity} instead: $L - y'\,\dfrac{\partial L}{\partial y'} = C$ (first integral).
    \item \textbf{If there is an integral constraint} $J[y]=\int_a^b G\,dx = \text{const}$, form the augmented Lagrangian $H = L + \lambda G$ (isoperimetric method) and apply the E-L equation to $H$.
    \item \textbf{Solve the resulting ODE.} Apply the boundary conditions to determine integration constants.
    \item \textbf{Determine the Lagrange multiplier $\lambda$} (if present) by substituting into the constraint equation.
    \item \textbf{Verify:} confirm BCs, constraint satisfaction, and optionally compute $I$ along the extremal.
\end{enumerate}
\end{notebox}

\subsection{\texorpdfstring{Minimize $I[y]=\int_0^1(y')^2\,dx$ with $\int_0^1 y\,dx=2$; $y(0)=0$, $y(1)=1$}{Isoperimetric problem: augment functional with Lagrange multiplier; E-L gives parabola}}

\begin{problembox}
Minimize $I[y] = \int_0^1 (y')^2 \, dx$ subject to $J[y] = \int_0^1 y \, dx = 2$
and $y(0) = 0, y(1) = 1$.
\end{problembox}

We form the augmented functional $H[y] = I[y] + \lambda J[y] = \int_0^1 ((y')^2 + \lambda y) \, dx$.
The Lagrangian is $L(x, y, y') = (y')^2 + \lambda y$.
The Euler-Lagrange equation is:
$$ \frac{\partial L}{\partial y} - \frac{d}{dx}\left(\frac{\partial L}{\partial y'}\right) = 0 $$
$$ \lambda - \frac{d}{dx}(2y') = 0 \implies 2y'' = \lambda \implies y'' = \frac{\lambda}{2} $$
Integrating twice with respect to $x$:
$$ y'(x) = \frac{\lambda}{2}x + C_1 $$
$$ y(x) = \frac{\lambda}{4}x^2 + C_1 x + C_2 $$

\textbf{Apply Boundary Conditions:}
1. $y(0) = 0 \implies C_2 = 0$.
2. $y(1) = 1 \implies \frac{\lambda}{4}(1)^2 + C_1(1) = 1 \implies C_1 = 1 - \frac{\lambda}{4}$.
So, $y(x) = \frac{\lambda}{4}x^2 + \left(1-\frac{\lambda}{4}\right)x$.

\textbf{Apply Integral Constraint:}
$$ \int_0^1 \left[\frac{\lambda}{4}x^2 + \left(1-\frac{\lambda}{4}\right)x\right] \, dx = 2 $$
$$ \left[ \frac{\lambda x^3}{12} + \frac{(1-\lambda/4)x^2}{2} \right]_0^1 = 2 $$
$$ \lambda\left(\frac{1}{12} - \frac{1}{8}\right) = 2 - \frac{1}{2} = \frac{3}{2} $$
$$ \lambda\left(\frac{2-3}{24}\right) = \frac{3}{2} $$
$$ \lambda\left(-\frac{1}{24}\right) = \frac{3}{2} \implies \lambda = -36 $$

Substituting $\lambda = -36$ back:
$$ y(x) = \frac{-36}{4}x^2 + \left(1 - \frac{-36}{4}\right)x = -9x^2 + 10x $$

\begin{summarybox}
The minimizing curve is:
$$ y(x) = -9x^2 + 10x $$
\end{summarybox}

\newpage
\subsection{\texorpdfstring{Extremize $I[y]=\int_0^2\!\left(y-\sqrt{1+y'^2}\,\right)dx$; $y(0)=y(2)=8$}{Apply Beltrami identity; ODE solves to circular arc $(x-1)^2+(y-8)^2=1$}}

\begin{problembox}
Extremize:
$$I[y]=\int_{0}^{2}\Big( y-\sqrt{1+y'^2}\Big)\,dx, \quad y(0)=y(2)=8$$
\end{problembox}

The Lagrangian is $F(y, y') = y - \sqrt{1+y'^2}$, which does not depend explicitly on $x$.

\textbf{Apply Beltrami Identity:}
$$F - y'\frac{\partial F}{\partial y'} = C$$

Compute $\displaystyle\frac{\partial F}{\partial y'} = -\frac{y'}{\sqrt{1+y'^2}}$. Substituting:
$$y - \sqrt{1+y'^2} + \frac{y'^2}{\sqrt{1+y'^2}} = C$$
$$y - \frac{1}{\sqrt{1+y'^2}} = C$$

\textbf{Solve the ODE:}
Rearrange: $\displaystyle\sqrt{1+y'^2} = \frac{1}{y-C}$

Squaring: $\displaystyle y'^2 = \frac{1-(y-C)^2}{(y-C)^2}$

Separating variables with substitution $y - C = \sin\theta$:
$$\sin\theta \, d\theta = \pm dx \implies \cos\theta = A \mp x$$

Thus: $\displaystyle(x-A)^2 + (y-C)^2 = 1$ (circle of radius 1)

\textbf{Apply Boundary Conditions:}
At $x=0$ and $x=2$: $y = 8$ gives
$$\sqrt{1-A^2} = \sqrt{1-(A-2)^2} \implies A = 1$$
$$8 = C + 0 \implies C = 8$$

\begin{summarybox}
$$(x-1)^2 + (y-8)^2 = 1$$
$$y(x) = 8 \pm \sqrt{1-(x-1)^2}$$

Upper arc: $y = 8 + \sqrt{1-(x-1)^2}$

Lower arc: $y = 8 - \sqrt{1-(x-1)^2}$
\end{summarybox}

\newpage
% ---------------------------------------------------------------------
% NEW PROBLEM (Section 5)
% ---------------------------------------------------------------------
\subsection{\texorpdfstring{Find extremal of $I[y]=\int_0^{\pi/2}(y'^2-y^2)\,dx$; $y(0)=0$, $y(\pi/2)=1$}{CoV: Euler-Lagrange gives $y''+y=0$; general solution is trigonometric; unique extremal $y=\sin x$}}

\begin{problembox}
Find the extremal of the functional:
$$I[y] = \int_0^{\pi/2}\!\left(y'^2 - y^2\right)dx$$
subject to the boundary conditions $y(0) = 0$, $y\!\left(\dfrac{\pi}{2}\right) = 1$.

Also compute the value of $I$ along the extremal.
\end{problembox}

\textbf{Step 1: Write down the Euler-Lagrange equation.}
The Lagrangian is $L(x, y, y') = y'^2 - y^2$. The Euler-Lagrange equation is:
\[
    \frac{\partial L}{\partial y} - \frac{d}{dx}\frac{\partial L}{\partial y'} = 0.
\]
Computing each term:
\[
    \frac{\partial L}{\partial y} = -2y, \qquad \frac{\partial L}{\partial y'} = 2y', \qquad \frac{d}{dx}(2y') = 2y''.
\]
The Euler-Lagrange equation becomes:
\[
    -2y - 2y'' = 0 \implies y'' + y = 0.
\]

\textbf{Step 2: Solve the ODE.}
The characteristic equation $r^2 + 1 = 0$ has roots $r = \pm i$, giving the general solution:
\[
    y(x) = A\cos x + B\sin x.
\]

\textbf{Step 3: Apply boundary conditions.}
\begin{align*}
    y(0) = 0: &\quad A\cos 0 + B\sin 0 = A = 0. \\
    y\!\left(\tfrac{\pi}{2}\right) = 1: &\quad B\sin\!\tfrac{\pi}{2} = B = 1.
\end{align*}

\begin{summarybox}
The unique extremal is:
$$y(x) = \sin x$$
\end{summarybox}

\textbf{Step 4: Compute the value of $I$ along the extremal.}
With $y = \sin x$ and $y' = \cos x$:
\[
    I[\sin x] = \int_0^{\pi/2}\!\left(\cos^2 x - \sin^2 x\right)dx = \int_0^{\pi/2}\cos(2x)\,dx = \left[\frac{\sin(2x)}{2}\right]_0^{\pi/2} = 0.
\]

\textbf{Remark:} The value $I = 0$ is consistent with $y = \sin x$ being a stationary point; however, competitors such as $\tilde y = \sin x + \epsilon\sin(3x)$ (which also vanish at $0$ and $\frac{\pi}{2}$, but with $\tilde y(\pi/2) = \sin(\pi/2) + \epsilon\sin(3\pi/2) = 1 - \epsilon$) show the extremal is a saddle point of the functional rather than a minimum.

\newpage
% ---------------------------------------------------------------------
% PROBLEM 7
% ---------------------------------------------------------------------
\section{Green's Functions and Sturm-Liouville Theory}

\begin{notebox}[title={General Procedure: Constructing Green's Functions}]
To construct the Green's function $G(x,x')$ satisfying $\hat{L}\,G = \delta(x-x')$ for a self-adjoint Sturm--Liouville operator $\hat{L} = -\frac{d}{dx}\!\left(p(x)\frac{d}{dx}\right)+q(x)$ with given BCs:
\begin{enumerate}
    \item \textbf{Solve the homogeneous equation} $\hat{L}y = 0$ for two linearly independent solutions.
    \item \textbf{Select basis functions:} Choose $y_1(x)$ to satisfy the \emph{left} BC and $y_2(x)$ to satisfy the \emph{right} BC.
    \item \textbf{Form the piecewise Green's function} (symmetric for self-adjoint operators):
    $G(x,x') = A\begin{cases} y_1(x)\,y_2(x') & x < x' \\ y_1(x')\,y_2(x) & x > x' \end{cases}$
    \item \textbf{Apply the jump condition} on the derivative at $x = x'$:
    $\left[\frac{\partial G}{\partial x}\right]_{x'^-}^{x'^+} = -\frac{1}{p(x')}$
    to determine the constant $A$.
    \item \textbf{Write the solution} to $\hat{L}u = f$ as $u(x) = \int G(x,x')\,f(x')\,dx'$.
    \item \textbf{Verify:} symmetry $G(x,x')=G(x',x)$; boundary conditions; continuity at $x=x'$; correct derivative jump.
\end{enumerate}
\end{notebox}

\subsection{\texorpdfstring{Construct $G(x,x')$ for $\hat{L}y=-\tfrac{d}{dx}\!\left((1+x)y'\right)$ on $(0,1)$; $y(0)=0$, $y'(1)=0$}{Green's function for SL operator: left/right basis solutions, jump condition gives $A=1$}}

\begin{problembox}
Construct $G(x, x')$ for $\hat{L}y = -\frac{d}{dx}\left( (1+x) \frac{dy}{dx} \right)$ on $0 < x < 1$.
BCs: $y(0) = 0, \quad y'(1) = 0$.
\end{problembox}

The Green's function satisfies $\hat{L}G(x, x') = \delta(x-x')$.
First, solve the homogeneous equation $((1+x)y')' = 0$:
$$ (1+x)y' = C \implies y' = \frac{C}{1+x} \implies y(x) = C \ln(1+x) + D $$

\textbf{Step 1: Construct Basis Functions}
Find $y_1(x)$ satisfying the left BC $y(0)=0$:
$C \ln(1) + D = 0 \implies D = 0$.
$$ y_1(x) = \ln(1+x) $$

Find $y_2(x)$ satisfying the right BC $y'(1)=0$:
$y'(x) = \frac{C}{1+x}$. At $x=1$, $\frac{C}{2} = 0 \implies C=0$.
So $y_2(x) = \text{constant}$. We choose:
$$ y_2(x) = 1 $$

\textbf{Step 2: Construct G}
$$ G(x, x') = \begin{cases} A y_1(x)y_2(x') & 0 \le x < x' \\ A y_1(x')y_2(x) & x' < x \le 1 \end{cases} $$
Note: Since the Sturm-Liouville operator $\hat{L}$ is self-adjoint (it can be written in the form $-\frac{d}{dx}(p\frac{d}{dx})$ with appropriate boundary conditions), the Green's function must be symmetric: $G(x,x') = G(x',x)$.
$G(x, x') = \begin{cases} A \ln(1+x) & x < x' \\ A \ln(1+x') & x > x' \end{cases}$

\textbf{Step 3: Jump Condition}
For operator $L = -\frac{d}{dx}(p(x) \frac{d}{dx})$, the jump in the derivative is:
$$ \left. \frac{\partial G}{\partial x} \right|_{x'+\epsilon} - \left. \frac{\partial G}{\partial x} \right|_{x'-\epsilon} = -\frac{1}{p(x')} $$
Here $p(x) = 1+x$.
Derivative for $x > x'$: $\frac{\partial}{\partial x}(A \ln(1+x')) = 0$.
Derivative for $x < x'$: $\frac{\partial}{\partial x}(A \ln(1+x)) = \frac{A}{1+x}$.
Jump at $x=x'$:
$$ 0 - \frac{A}{1+x'} = -\frac{1}{1+x'} \implies A = 1 $$

\begin{summarybox}
Let $x_<$ be the smaller of $x, x'$ and $x_>$ be the larger.
$$ G(x, x') = \ln(1 + x_<) $$
\end{summarybox}

\newpage

\subsection{\texorpdfstring{$G'' + k^2 G = \delta(x-x_0)$ on $\mathbb{R}$; outgoing-wave BCs $G\sim e^{\pm ikx}$ as $x\to\pm\infty$}{Green's function for 1D Helmholtz on real line; continuity and jump condition give $G=\frac{e^{ik|x-x_0|}}{2ik}$}}

\begin{problembox}
Consider the Sturm-Liouville operator $L = \frac{d^2}{dx^2} + k^2$ on $-\infty < x < \infty$.

 Find the Green's function satisfying $LG(x, x_0) = \delta(x - x_0)$ with outgoing boundary conditions ($G \sim e^{\pm ikx}$ as $x \to \pm\infty$).

\end{problembox}
We solve the differential equation:
$$ \left( \frac{d^2}{dx^2} + k^2 \right) G(x, x_0) = \delta(x - x_0) $$

For $x \neq x_0$, the homogeneous equation $u'' + k^2 u = 0$ admits solutions $e^{\pm ikx}$.

Applying radiation boundary conditions (outgoing waves from source $x_0$):
$$ G(x, x_0) = \begin{cases} A e^{-ikx} & x < x_0 \\ B e^{ikx} & x > x_0 \end{cases} $$

\textbf{Step 1: Continuity at $x = x_0$.}
$$ A e^{-ikx_0} = B e^{ikx_0} \implies A = B e^{2ikx_0} $$

\textbf{Step 2: Jump condition (integrate ODE across $x_0$).}
$$ \left.\frac{dG}{dx}\right|_{x_0^+} - \left.\frac{dG}{dx}\right|_{x_0^-} = 1 $$
Substituting the derivatives:
$$ (ik B e^{ikx_0}) - (-ik A e^{-ikx_0}) = 1 $$
Using $A = B e^{2ikx_0}$:
$$ 2ik B e^{ikx_0} = 1 \implies B = \frac{1}{2ik} e^{-ikx_0} $$

Consequently, $A = \frac{1}{2ik} e^{ikx_0}$.

\begin{summarybox}[title={Part A Answer}]
$$ G(x, x_0) = \frac{1}{2ik} e^{ik|x - x_0|} $$
\end{summarybox}



\newpage

\subsection{\texorpdfstring{$\hat{L}u=\tfrac{d}{dx}(e^{-x}u')$ on $(0,\ln 2)$: (A) find $G(x,x_0)$ with $G(0)=G(\ln 2)=0$; (B) solve $\hat{L}u=-\lambda u$, $u(0)=1$, $u(\ln 2)=2$ to first order in $\lambda$}{(A) Green's function for exponential-weight SL operator; (B) perturbation expansion: zeroth order $u_0=e^x$}}

\begin{problembox}
Given the Sturm-Liouville operator $\hat{L}u = \frac{d}{dx} \left( e^{-x} \frac{du}{dx} \right)$ on $0 < x < \ln 2$.

(A) Find the Green's function $G(x, x_0)$ satisfying $\hat{L}G = -\delta(x-x_0)$ with Dirichlet BCs $G(0)=G(\ln 2)=0$.

(B) Solve $\hat{L}u = -\lambda u$ with $u(0)=1, u(\ln 2)=2$ to first order in $\lambda$.
\end{problembox}

\textbf{(A) Green's Function Construction}

The homogeneous equation $(e^{-x}u')' = 0$ yields solutions of the form $u(x) = C_1 e^x + C_2$.
We construct solutions satisfying the boundary conditions:
\begin{itemize}
    \item \textbf{Left ($x < x_0$):} Satisfies $u(0)=0$.
    $$ u_L(x) = e^x - 1 $$
    \item \textbf{Right ($x > x_0$):} Satisfies $u(\ln 2)=0$.
    $$ u_R(x) = e^x - 2 $$
\end{itemize}

The Green's function is continuous at $x_0$:
$$ G(x, x_0) = A \begin{cases} (e^x - 1)(e^{x_0} - 2) & 0 \le x \le x_0 \\ (e^{x_0} - 1)(e^x - 2) & x_0 \le x \le \ln 2 \end{cases} $$

Apply the jump condition on the derivative: $[e^{-x} G']_{x_0^-}^{x_0^+} = -1$.
$$ e^{-x_0} \left[ G'(x_0^+) - G'(x_0^-) \right] = -1 $$
$$ e^{-x_0} \left[ A(e^{x_0}-1)e^{x_0} - A(e^{x_0}-2)e^{x_0} \right] = -1 $$
$$ A [ (e^{x_0}-1) - (e^{x_0}-2) ] = -1 \implies A(1) = -1 \implies A = -1 $$

\begin{summarybox}[title={Part A Answer}]
$$ G(x, x_0) = \begin{cases} (e^x - 1)(2 - e^{x_0}) & 0 \le x \le x_0 \\ (e^{x_0} - 1)(2 - e^x) & x_0 \le x \le \ln 2 \end{cases} $$
\end{summarybox}

\textbf{(B) Perturbation Expansion}

We seek a solution $u(x) \approx u_0(x) + \lambda u_1(x)$.
Substituting into $\hat{L}u = -\lambda u$:
$$ \hat{L}(u_0 + \lambda u_1) = -\lambda(u_0 + \lambda u_1) \approx -\lambda u_0 $$

\textbf{Order $\lambda^0$ (Zeroth Order):}
$\hat{L}u_0 = 0$ with $u_0(0)=1, u_0(\ln 2)=2$.
$$ u_0(x) = C_1 e^x + C_2 $$
Using BCs: $C_1 + C_2 = 1$ and $2C_1 + C_2 = 2 \implies C_1 = 1, C_2 = 0$.
$$ u_0(x) = e^x $$

\textbf{Order $\lambda^1$ (First Order):}
$\hat{L}u_1 = -u_0 = -e^x$ with homogeneous BCs $u_1(0)=0, u_1(\ln 2)=0$.
$$ (e^{-x} u_1')' = -e^x \implies e^{-x} u_1' = -e^x + A \implies u_1' = -e^{2x} + Ae^x $$
Integrating again:
$$ u_1(x) = -\frac{1}{2}e^{2x} + Ae^x + B $$
Apply BCs:
\begin{align*}
    u_1(0) &= -0.5 + A + B = 0 \\
    u_1(\ln 2) &= -0.5(4) + 2A + B = 0 \implies 2A + B = 2
\end{align*}
Subtracting the first from the second: $(2A+B) - (A+B) = 2 - 0.5 \implies A = 1.5$.
Then $B = 0.5 - 1.5 = -1$.

$$ u_1(x) = -\frac{1}{2}e^{2x} + \frac{3}{2}e^x - 1 $$

\begin{summarybox}[title={Part B Answer}]
$$ u(x) \approx e^x + \lambda \left( -\frac{1}{2}e^{2x} + \frac{3}{2}e^x - 1 \right) $$
\end{summarybox}

\newpage
% ---------------------------------------------------------------------
% NEW PROBLEM (Section 6)
% ---------------------------------------------------------------------
\subsection{\texorpdfstring{Construct $G(x,x')$ for $-y''=f$ on $(0,1)$; $y(0)=y(1)=0$}{Standard 1D Green's function: piecewise linear; $G(x,x')=x_<(1-x_>)$}}

\begin{problembox}
Construct the Green's function $G(x, x')$ for the boundary-value problem:
$$-y''(x) = f(x), \quad 0 < x < 1, \qquad y(0) = 0, \quad y(1) = 0.$$
That is, find $G(x,x')$ satisfying $-G_{xx} = \delta(x-x')$ with the same boundary conditions.
Use $G$ to write the formal solution for general $f$.
\end{problembox}

\textbf{Step 1: Homogeneous solutions satisfying each BC.}

The homogeneous equation $y'' = 0$ has general solution $y = Ax + B$.

\begin{itemize}
    \item \textbf{Left solution $y_1$} (satisfies $y_1(0)=0$): $B=0$, so $y_1(x) = x$.
    \item \textbf{Right solution $y_2$} (satisfies $y_2(1)=0$): $A + B = 0 \Rightarrow B = -A$, so $y_2(x) = 1-x$ (choosing $A=-1$).
\end{itemize}

\textbf{Step 2: Write the piecewise Green's function.}

The standard form for a self-adjoint operator $Ly = -y''$ is:
\[
    G(x,x') = \frac{1}{p(x') W(x')} \begin{cases} y_1(x)\,y_2(x') & x < x' \\ y_1(x')\,y_2(x) & x > x' \end{cases}
\]
where $p(x) = 1$ (coefficient of $y''$) and $W(x') = y_1 y_2' - y_1' y_2$.

\textbf{Step 3: Compute the Wronskian.}
\[
    W(x') = y_1(x') y_2'(x') - y_1'(x') y_2(x')
    = x'(-1) - (1)(1-x') = -x' - 1 + x' = -1.
\]

\textbf{Step 4: Determine prefactor from jump condition.}
For operator $L = -\frac{d^2}{dx^2}$, the jump condition is:
\[
    G'(x'^+,x') - G'(x'^-,x') = -\frac{1}{p(x')} = -1.
\]
With the Wronskian $W = -1$ and $p = 1$:
\[
    \frac{1}{p\,W} = \frac{1}{1\cdot(-1)} = -1,
\]
so:
\[
    G(x,x') = (-1)\begin{cases} x(1-x') & x < x' \\ x'(1-x) & x > x' \end{cases}
    = \begin{cases} x(1-x') & x < x' \\ x'(1-x) & x > x' \end{cases}.
\]
(The sign works out correctly since $W = -1$ and the operator is $-\partial_{xx}$.)

\begin{summarybox}
$$G(x, x') = x_<\,(1 - x_>)$$
where $x_< = \min(x,x')$ and $x_> = \max(x,x')$.

The solution of $-y'' = f$ with $y(0)=y(1)=0$ is:
$$y(x) = \int_0^1 G(x, x') f(x')\,dx' = (1-x)\int_0^x x' f(x')\,dx' + x\int_x^1 (1-x') f(x')\,dx'$$
\end{summarybox}

\textbf{Verification.}
\begin{itemize}
    \item \textbf{Symmetry:} $G(x,x') = G(x',x)$ (self-adjoint operator). \checkmark
    \item \textbf{BCs:} $G(0,x') = 0 \cdot (1-x') = 0$ and $G(1,x') = x'(1-1) = 0$. \checkmark
    \item \textbf{Continuity at $x=x'$:} $x'(1-x') = x'(1-x')$. \checkmark
    \item \textbf{Jump in derivative:} For $x < x'$, $\partial_x G = 1-x'$; for $x > x'$, $\partial_x G = -x'$. Jump $= -x' - (1-x') = -1 = -1/p(x')$. \checkmark
\end{itemize}

\newpage
% ---------------------------------------------------------------------
% PROBLEM 8
% ---------------------------------------------------------------------
\section{Integral Equations and Generalized Functions}
\vspace{-0.5cm}
\begin{notebox}[title={General Procedure: Integral Equations \& Dirac Delta}]
\textbf{Fredholm Equations with Separable Kernel} $K(x,t)=\sum_i \phi_i(x)\psi_i(t)$:
\begin{enumerate}
    \item Define scalar unknowns $C_i = \int \psi_i(t)\,\varphi(t)\,dt$.
    \item Express $\varphi(x)$ in terms of the $C_i$ and $\phi_i(x)$.
    \item Substitute back into the definitions of $C_i$ to obtain a linear system $M\mathbf{c} = \mathbf{b}$.
    \item Solve for $C_i$; a unique solution exists when $\det(M)\ne 0$.
    \item \textbf{Eigenvalues:} set $\det(M)=0$ to find the values of $\lambda$ where no unique solution exists.
\end{enumerate}
\textbf{Volterra Equations with Convolution Kernel:}
\begin{enumerate}
    \item Simplify by multiplying out the convolution factor (e.g.\ multiply both sides by $e^{\alpha x}$).
    \item Differentiate both sides to convert the integral equation into an ODE.
    \item Solve the ODE using standard methods (integrating factor, etc.).
    \item Verify by substituting back into the original integral equation.
\end{enumerate}
\textbf{Dirac Delta Computations:}
\begin{itemize}
    \item \textbf{Composition rule:} If $g(x)$ has simple roots $x_i$, then $\delta(g(x))=\sum_i \frac{\delta(x-x_i)}{|g'(x_i)|}$.
    \item \textbf{Sifting property:} $\int f(x)\,\delta(x-a)\,dx = f(a)$.
    \item \textbf{Distributional derivatives:} A jump $J$ in $f^{(n-1)}$ at $x_0$ produces a term $J\,\delta(x-x_0)$ in $f^{(n)}$.
\end{itemize}
\end{notebox}

\subsection{\texorpdfstring{Solve $\varphi(x)=x+\lambda\int_0^1(xt^2+x^2t)\varphi(t)\,dt$; find eigenvalues $\lambda$}{Fredholm IE with separable kernel; reduce to $2\times 2$ linear system; eigenvalues from $\det M=0$}}

\begin{problembox}
Solve for $\varphi(x)$:
$$\varphi(x) = x + \lambda \int_0^1 (x t^2 + x^2 t) \varphi(t) \, dt$$
Determine eigenvalues $\lambda$.
\end{problembox}

This is a Fredholm equation with a separable kernel $K(x,t) = x t^2 + x^2 t$.
We can rewrite the equation as:
$$ \varphi(x) = x + \lambda x \int_0^1 t^2 \varphi(t) \, dt + \lambda x^2 \int_0^1 t \varphi(t) \, dt $$
Define the constants:
$$ C_1 = \int_0^1 t^2 \varphi(t) \, dt, \quad C_2 = \int_0^1 t \varphi(t) \, dt $$
Then the solution form is:
$$ \varphi(x) = x + \lambda C_1 x + \lambda C_2 x^2 = x(1+\lambda C_1) + \lambda C_2 x^2 $$
Substitute this back into the definitions of $C_1$ and $C_2$:

$$ C_1 = \int_0^1 t^2 [t(1+\lambda C_1) + \lambda C_2 t^2] \, dt = \int_0^1 [(1+\lambda C_1)t^3 + \lambda C_2 t^4] \, dt $$
$$ C_1 = (1+\lambda C_1)\frac{1}{4} + \lambda C_2 \frac{1}{5} $$
$$ C_1 \left(1 - \frac{\lambda}{4}\right) - \frac{\lambda}{5}C_2 = \frac{1}{4} \quad \dots (1) $$

$$ C_2 = \int_0^1 t [t(1+\lambda C_1) + \lambda C_2 t^2] \, dt = \int_0^1 [(1+\lambda C_1)t^2 + \lambda C_2 t^3] \, dt $$
$$ C_2 = (1+\lambda C_1)\frac{1}{3} + \lambda C_2 \frac{1}{4} $$
$$ - \frac{\lambda}{3}C_1 + C_2 \left(1 - \frac{\lambda}{4}\right) = \frac{1}{3} \quad \dots (2) $$

This is a linear system $M \mathbf{c} = \mathbf{b}$. A unique solution exists if $\det(M) \neq 0$.
$$ D(\lambda) = \left(1 - \frac{\lambda}{4}\right)^2 - \left(-\frac{\lambda}{5}\right)\left(-\frac{\lambda}{3}\right) = 1 - \frac{\lambda}{2} + \frac{\lambda^2}{16} - \frac{\lambda^2}{15} $$
$$ D(\lambda) = 1 - \frac{\lambda}{2} + \lambda^2 \frac{15-16}{240} = 1 - \frac{\lambda}{2} - \frac{\lambda^2}{240} $$

The equation has a unique solution when $D(\lambda) \neq 0$. Setting $D(\lambda) = 0$ gives the eigenvalues:
$$ 1 - \frac{\lambda}{2} - \frac{\lambda^2}{240} = 0 \implies \lambda = -60 \pm 4\sqrt{60}$$


For $\lambda$ not equal to these eigenvalues, we can solve the linear system (1)-(2) for $C_1, C_2$ using Cramer's rule:
$$ C_1 = \frac{\frac{1}{4}(1-\frac{\lambda}{4}) - \frac{1}{3}(-\frac{\lambda}{5})}{D(\lambda)}, \quad C_2 = \frac{(1-\frac{\lambda}{4})\frac{1}{3} - \frac{1}{4}(-\frac{\lambda}{3})}{D(\lambda)} $$

The general solution is then:
$$ \varphi(x) = (1+\lambda C_1)x + \lambda C_2 x^2 $$

\newpage

% ---------------------------------------------------------------------
% PROBLEM 9
% ---------------------------------------------------------------------
\subsection{\texorpdfstring{Evaluate $I=\int_{-\infty}^{\infty}(x^3+2x)\,\delta(x^2-3x+2)\,dx$}{Apply delta composition rule $\delta(g(x))=\sum_i\delta(x-x_i)/|g'(x_i)|$; roots at $x=1,2$; $I=15$}}

\begin{problembox}
Evaluate the integral:
$$I = \int_{-\infty}^{\infty} (x^3 + 2x) \, \delta(x^2 - 3x + 2) \, dx$$
\end{problembox}


We use the composition rule for the Dirac delta function. If $g(x)$ has simple roots $x_i$, then:
$$ \delta(g(x)) = \sum_i \frac{\delta(x - x_i)}{|g'(x_i)|} $$

\textbf{Step 1: Find roots of the argument}
Let $g(x) = x^2 - 3x + 2$.
Roots: $(x-1)(x-2) = 0 \implies x_1 = 1, x_2 = 2$.

\textbf{Step 2: Evaluate derivative at roots}
$g'(x) = 2x - 3$.
At $x_1 = 1$: $g'(1) = 2(1) - 3 = -1 \implies |g'(1)| = 1$.
At $x_2 = 2$: $g'(2) = 2(2) - 3 = 1 \implies |g'(2)| = 1$.

\textbf{Step 3: Expand the Delta Function}
$$ \delta(x^2 - 3x + 2) = \frac{\delta(x-1)}{1} + \frac{\delta(x-2)}{1} = \delta(x-1) + \delta(x-2) $$

\textbf{Step 4: Evaluate the Integral}
Let $f(x) = x^3 + 2x$.
$$ I = \int_{-\infty}^{\infty} f(x) [\delta(x-1) + \delta(x-2)] \, dx $$
Using the sifting property $\int f(x)\delta(x-a)dx = f(a)$:
$$ I = f(1) + f(2) $$
Calculate $f(1)$:
$$ f(1) = 1^3 + 2(1) = 3 $$
Calculate $f(2)$:
$$ f(2) = 2^3 + 2(2) = 8 + 4 = 12 $$

$$ I = 3 + 12 = 15 $$

\begin{summarybox}
$$ I = 15 $$
\end{summarybox}

\newpage

% ---------------------------------------------------------------------
% PROBLEM 10
% ---------------------------------------------------------------------
\subsection{\texorpdfstring{(A) Evaluate $\int_{-\infty}^{\infty}\delta(e^{|x|}\sin x)\,dx$; (B) distributional derivatives of $f(x)=1+e^{-|x|}$}{(A) Roots at $x=n\pi$; geometric series gives $\coth(\pi/2)$; (B) $f''=e^{-|x|}-2\delta(x)$}}

\begin{problembox}
\textbf{Part A:} Calculate the integral:
$\int_{-\infty}^{\infty} \delta(e^{|x|} \sin x) \, dx$

\textbf{Part B:} Calculate the distributional derivative of the function $f(x) = 1 + e^{-|x|}$.
\end{problembox}

\textbf{Part A: Delta Function with Composite Argument}

We use the composition rule for the Dirac delta function. Let $g(x) = e^{|x|} \sin x$.
The support of the delta function $\delta(g(x))$ is the set of roots $x_n$ where $g(x_n) = 0$.
Since $e^{|x|} > 0$, the roots are determined by $\sin x = 0$:
$$ x_n = n\pi, \quad n \in \mathbb{Z} $$

The derivative is $g'(x) = \text{sgn}(x) e^{|x|} \sin x + e^{|x|} \cos x$.
At the roots $x_n = n\pi$, the sine term vanishes:
$$ g'(n\pi) = e^{|n\pi|} \cos(n\pi) = (-1)^n e^{|n|\pi} $$
$$ |g'(n\pi)| = e^{|n|\pi} $$

Using the expansion $\delta(g(x)) = \sum_n \frac{\delta(x - x_n)}{|g'(x_n)|}$:
$$ \delta(e^{|x|} \sin x) = \sum_{n=-\infty}^{\infty} \frac{\delta(x - n\pi)}{e^{|n|\pi}} $$

The integral is the sum of the weights:
$$ I = \sum_{n=-\infty}^{\infty} e^{-|n|\pi} = 1 + 2\sum_{n=1}^{\infty} (e^{-\pi})^n $$

This is a geometric series with $r = e^{-\pi} < 1$:
$$ \sum_{n=1}^{\infty} r^n = \frac{r}{1-r} = \frac{e^{-\pi}}{1 - e^{-\pi}} = \frac{1}{e^{\pi} - 1} $$
$$ I = 1 + \frac{2}{e^{\pi} - 1} = \frac{e^{\pi} - 1 + 2}{e^{\pi} - 1} = \frac{e^{\pi} + 1}{e^{\pi} - 1} $$

Recalling that $\coth(x) = \frac{e^{2x} + 1}{e^{2x} - 1}$, we can verify:
$$ \coth\left(\frac{\pi}{2}\right) = \frac{e^{\pi} + 1}{e^{\pi} - 1} $$

Therefore, $I = \coth\left(\frac{\pi}{2}\right)$.

\begin{summarybox}[title={Part A Answer}]
$$ \int_{-\infty}^{\infty} \delta(e^{|x|} \sin x) \, dx = \coth\left(\frac{\pi}{2}\right) $$
\end{summarybox}

\textbf{Part B: Distributional Derivative}

Given $f(x) = 1 + e^{-|x|}$:
$$ f(x) = \begin{cases} 1 + e^x & x < 0 \\ 1 + e^{-x} & x > 0 \end{cases} $$

1. \textbf{First Derivative:}
The function $f(x)$ is continuous at $x=0$ since $\lim_{x\to 0^-} f(x) = 2$ and $\lim_{x\to 0^+} f(x) = 2$.
Therefore, the distributional derivative involves no delta functions and equals the classical derivative defined piecewise:
$$ f'(x) = \begin{cases} e^x & x < 0 \\ -e^{-x} & x > 0 \end{cases} $$

This can be written compactly as $f'(x) = -\text{sgn}(x)e^{-|x|}$ for $x \neq 0$, where $\text{sgn}(x)$ is the sign function.

2. \textbf{Second Derivative:}
Check for discontinuity in $f'(x)$ at $x=0$:
$$ \lim_{x\to 0^-} f'(x) = 1, \quad \lim_{x\to 0^+} f'(x) = -1 $$
Jump $J = -1 - 1 = -2$.
The second derivative contains a delta function proportional to the jump:
$$ f''(x) = \{f''(x)\}_{classical} + J\delta(x) $$
$$ \{f''(x)\}_{classical} = \begin{cases} e^x & x < 0 \\ e^{-x} & x > 0 \end{cases} = e^{-|x|} $$
$$ f''(x) = e^{-|x|} - 2\delta(x) $$

\begin{summarybox}[title={Part B Answer}]
$$ f'(x) = -\text{sgn}(x) e^{-|x|} $$
(Standard distributions; no delta part)
\end{summarybox}

\newpage

% ---------------------------------------------------------------------
% PROBLEM 11
% ---------------------------------------------------------------------
\subsection{\texorpdfstring{Find eigenvalue $\lambda$ for $f(\mathbf{x})=\lambda\iint_{|\mathbf{x}'|<R}\frac{f(\mathbf{x}')}{1+3|\mathbf{x}'|}\,d\mathbf{x}'$ in disk $|\mathbf{x}'|<R$}{Kernel independent of $\mathbf{x}$ implies $f=\mathrm{const}$; eigenvalue from 2D polar integral over disk}}

\begin{problembox}
Find the eigenvalue $\lambda$ and eigenfunction $f(\mathbf{x})$ for the integral equation:
$$f(\mathbf{x}) = \lambda \iint_{|\mathbf{x}'|<R} \frac{f(\mathbf{x}')}{1+3|\mathbf{x}'|} d\mathbf{x}'$$
\end{problembox}

\textbf{Solution:}

The kernel $K(\mathbf{x}, \mathbf{x}') = \frac{1}{1+3|\mathbf{x}'|}$ is independent of $\mathbf{x}$. Therefore, the right-hand side is a constant (independent of $\mathbf{x}$).
This implies that the eigenfunction must be a constant:
$$ f(\mathbf{x}) = C $$

Substituting $f(\mathbf{x}) = C$ into the equation:
$$ C = \lambda \iint_{|\mathbf{x}'|<R} \frac{C}{1+3|\mathbf{x}'|} d\mathbf{x}' $$
Assuming $C \neq 0$, we divide by $C$:
$$ 1 = \lambda \iint_{|\mathbf{x}'|<R} \frac{1}{1+3|\mathbf{x}'|} d\mathbf{x}' $$

We evaluate the integral over the disk $|\mathbf{x}'|<R$ using 2D polar coordinates. Setting $\mathbf{x}' = (r\cos\theta, r\sin\theta)$ with $|\mathbf{x}'| = r$ and area element $d\mathbf{x}' = r\, dr\, d\theta$:
$$ I = \int_0^{2\pi} d\theta \int_0^R \frac{r}{1+3r} dr = 2\pi \int_0^R \frac{r}{1+3r} dr $$

Use substitution $u = 1+3r \implies r = \frac{u-1}{3}, \, dr = \frac{du}{3}$.
Limits: $r=0 \to u=1$; $r=R \to u=1+3R$.
$$ \int_0^R \frac{r}{1+3r} dr = \int_1^{1+3R} \frac{u-1}{3u} \frac{du}{3} = \frac{1}{9} \int_1^{1+3R} \left(1 - \frac{1}{u}\right) du $$
$$ = \frac{1}{9} \Big[ u - \ln|u| \Big]_1^{1+3R} $$
$$ = \frac{1}{9} (3R - \ln(1+3R)) $$

So the condition for $\lambda$ is:
$$ 1 = \lambda \cdot 2\pi \cdot \frac{1}{9} (3R - \ln(1+3R)) $$

\begin{summarybox}
\textbf{Eigenfunction:} $f(\mathbf{x}) = \text{const}$

\textbf{Eigenvalue:}
$$ \lambda = \frac{9}{2\pi (3R - \ln(1+3R))} $$
\end{summarybox}

\newpage

% ---------------------------------------------------------------------
% PROBLEM 12
% ---------------------------------------------------------------------

\newpage

% ---------------------------------------------------------------------
% PROBLEM 12
% ---------------------------------------------------------------------
\subsection{\texorpdfstring{Solve $f(x)=\theta(x)e^{-\alpha x}+\lambda\int_{-\infty}^{x}e^{-\alpha(x-y)}f(y)\,dy$, $\alpha>\lambda>0$}{Volterra convolution IE: multiply by $e^{\alpha x}$, differentiate to get $f'+( \alpha-\lambda)f=\delta(x)$}}

\begin{problembox}
Solve the Volterra integral equation:
$$f(x) = \theta(x)e^{-\alpha x} + \lambda \int_{-\infty}^{x} e^{-\alpha(x-y)} f(y) dy$$
where $\theta(x)$ is the Heaviside step function and $\alpha > \lambda > 0$.
\end{problembox}

\textbf{Strategy:} This is a Volterra integral equation of the second kind with a convolution-type kernel. The key insight is to convert it to an Ordinary Differential Equation (ODE) by differentiation.

\textbf{Step 1: Simplify using exponential properties}

Multiply both sides by $e^{\alpha x}$ to factor the exponential term in the kernel:
\begin{equation}
    e^{\alpha x} f(x) = \theta(x) + \lambda \int_{-\infty}^{x} e^{\alpha y} f(y) dy
\end{equation}

The right-hand side now has a simpler structure: a Heaviside function plus an integral with variable upper limit.

\textbf{Step 2: Differentiate to obtain an ODE}

Apply $\frac{d}{dx}$ to both sides. For the left side, use the product rule. For the right side, use:
\begin{itemize}
    \item Distributional derivative: $\theta'(x) = \delta(x)$
    \item Leibniz rule for differentiation under the integral: $\frac{d}{dx}\int_{-\infty}^{x} g(y) dy = g(x)$
\end{itemize}

\begin{align}
    \frac{d}{dx} \left( e^{\alpha x} f(x) \right) &= \delta(x) + \lambda e^{\alpha x} f(x) \\
    \alpha e^{\alpha x} f(x) + e^{\alpha x} f'(x) &= \delta(x) + \lambda e^{\alpha x} f(x)
\end{align}

Dividing through by $e^{\alpha x}$ (noting $e^{-\alpha x}\delta(x) = \delta(x)$ for $x=0$):
\begin{equation}
    f'(x) + (\alpha - \lambda)f(x) = \delta(x)
    \label{eq:volterra_ode}
\end{equation}

\begin{notebox}
Equation (\ref{eq:volterra_ode}) is a first-order inhomogeneous linear ODE with a Dirac delta forcing term. This is equivalent to finding the Green's function for the operator $\mathcal{L} = \frac{d}{dx} + (\alpha - \lambda)$.
\end{notebox}

\textbf{Step 3: Solve using integrating factor}

The integrating factor is $I(x) = e^{(\alpha-\lambda)x}$. Multiply both sides by $I(x)$:
\begin{equation}
    e^{(\alpha-\lambda)x} f'(x) + (\alpha-\lambda) e^{(\alpha-\lambda)x} f(x) = e^{(\alpha-\lambda)x} \delta(x) = \delta(x)
\end{equation}

The left side is an exact differential:
\begin{equation}
    \frac{d}{dx} \left[ f(x) e^{(\alpha-\lambda)x} \right] = \delta(x)
\end{equation}

\textbf{Step 4: Integrate both sides}

Integrate from $-\infty$ to $x$. The integral of $\delta(y)$ over this interval is the Heaviside function:
\begin{equation}
    \int_{-\infty}^{x} \delta(y) dy = \theta(x)
\end{equation}

Since $f(x)$ must vanish as $x \to -\infty$ (from the original equation), the boundary term at $-\infty$ is zero:
\begin{equation}
    f(x) e^{(\alpha-\lambda)x} \Big|_{-\infty}^{x} = \theta(x)
\end{equation}

Therefore:
\begin{equation}
    f(x) e^{(\alpha-\lambda)x} = \theta(x)
\end{equation}

\textbf{Step 5: Solve for $f(x)$}

Multiply both sides by $e^{-(\alpha-\lambda)x}$:

\begin{summarybox}
$$ f(x) = \theta(x) e^{-(\alpha - \lambda)x} $$
\end{summarybox}

\textbf{Verification:}

\textbf{1. Boundary behavior:} For $x < 0$, both $\theta(x)=0$ and the integral limit make $f(x)=0$. \checkmark

\textbf{2. Exponential decay:} Since $\alpha > \lambda$, the exponent $-(\alpha-\lambda) < 0$, ensuring $f(x) \to 0$ as $x \to +\infty$. \checkmark

\textbf{3. Substitution check:} Substitute $f(x) = \theta(x) e^{-(\alpha-\lambda)x}$ into the original equation:
\begin{align*}
    \text{RHS} &= \theta(x)e^{-\alpha x} + \lambda \int_{-\infty}^{x} e^{-\alpha(x-y)} \theta(y) e^{-(\alpha-\lambda)y} dy \\
    &= \theta(x)e^{-\alpha x} + \lambda e^{-\alpha x} \int_{0}^{x} e^{(\alpha-\lambda)y} e^{-(\alpha-\lambda)y} dy \\
    &= \theta(x)e^{-\alpha x} + \lambda e^{-\alpha x} \cdot x = \theta(x)e^{-\alpha x} \left(1 + \lambda x\right) \quad \text{(for $x>0$)}
\end{align*}

Actually, let me recalculate. For $x>0$:
\begin{align*}
    \int_{0}^{x} e^{\alpha y} e^{-(\alpha-\lambda)y} dy &= \int_{0}^{x} e^{\lambda y} dy = \frac{1}{\lambda}(e^{\lambda x} - 1)
\end{align*}
This gives $f(x) = \theta(x)e^{-\alpha x}(1 + e^{\lambda x} - 1) = \theta(x)e^{-(\alpha-\lambda)x}$. \checkmark

\newpage
% ---------------------------------------------------------------------
% NEW PROBLEM (Section 7)
% ---------------------------------------------------------------------
\subsection{\texorpdfstring{Solve $\varphi(x)=\cos(\pi x)+\lambda\int_0^1\cos(\pi x)\cos(\pi t)\,\varphi(t)\,dt$; find eigenvalue $\lambda=2$}{Fredholm IE with rank-1 kernel; unique solution $\frac{2}{2-\lambda}\cos(\pi x)$ for $\lambda\ne 2$; $\lambda=2$ is eigenvalue with eigenfunction $\cos(\pi x)$}}

\begin{problembox}
Solve the Fredholm integral equation of the second kind:
$$\varphi(x) = \cos(\pi x) + \lambda \int_0^1 \cos(\pi x)\cos(\pi t)\,\varphi(t)\,dt$$

(a) Find the solution for $\lambda \neq 2$.

(b) Identify the eigenvalue and corresponding eigenfunction of the integral operator.
\end{problembox}

\textbf{Step 1: Exploit the separable (rank-1) kernel.}

The kernel $K(x,t) = \cos(\pi x)\cos(\pi t)$ factors as a product $\phi(x)\psi(t)$ with $\phi = \psi = \cos(\pi\cdot)$.
Define the scalar:
\[
    C = \int_0^1 \cos(\pi t)\,\varphi(t)\,dt.
\]
The integral equation becomes:
\[
    \varphi(x) = \cos(\pi x) + \lambda C \cos(\pi x) = (1 + \lambda C)\cos(\pi x).
\]

\textbf{Step 2: Self-consistent equation for $C$.}

Substitute $\varphi(t) = (1+\lambda C)\cos(\pi t)$ into the definition of $C$:
\[
    C = \int_0^1 \cos(\pi t)\cdot(1+\lambda C)\cos(\pi t)\,dt = (1+\lambda C)\int_0^1\cos^2(\pi t)\,dt.
\]
Evaluate the integral:
\[
    \int_0^1\cos^2(\pi t)\,dt = \frac{1}{2} + \frac{1}{2}\underbrace{\int_0^1\cos(2\pi t)\,dt}_{=0} = \frac{1}{2}.
\]
Therefore:
\[
    C = \frac{1+\lambda C}{2} \implies 2C = 1 + \lambda C \implies C(2-\lambda) = 1.
\]

\textbf{Step 3: Resolve the two cases.}

\textbf{(a) $\lambda \neq 2$:}
\[
    C = \frac{1}{2-\lambda}, \qquad
    \varphi(x) = \left(1 + \frac{\lambda}{2-\lambda}\right)\cos(\pi x) = \frac{2}{2-\lambda}\cos(\pi x).
\]

\textbf{(b) $\lambda = 2$:}
The equation $C(2-2) = 1$ reduces to $0 = 1$, a contradiction. Hence no solution exists. This means $\lambda = 2$ is an eigenvalue: the homogeneous equation
\[
    \psi(x) = 2\int_0^1\cos(\pi x)\cos(\pi t)\,\psi(t)\,dt
\]
is satisfied by $\psi(x) = \cos(\pi x)$.

\begin{summarybox}
\textbf{(a)} For $\lambda \neq 2$:
$$\varphi(x) = \frac{2}{2-\lambda}\cos(\pi x)$$
\textbf{(b)} The eigenvalue of the kernel $K(x,t) = \cos(\pi x)\cos(\pi t)$ on $[0,1]$ is $\lambda = 2$, with eigenfunction $\cos(\pi x)$.
\end{summarybox}

\textbf{Verification:} For $\lambda \neq 2$, substitute into the RHS:
\begin{align*}
    \text{RHS} &= \cos(\pi x) + \lambda\cos(\pi x)\int_0^1\cos(\pi t)\cdot\frac{2}{2-\lambda}\cos(\pi t)\,dt \\
    &= \cos(\pi x)\!\left(1 + \frac{2\lambda}{2-\lambda}\cdot\frac{1}{2}\right)
     = \cos(\pi x)\cdot\frac{2-\lambda+\lambda}{2-\lambda} = \frac{2}{2-\lambda}\cos(\pi x) = \varphi(x). \;\checkmark
\end{align*}

\end{document}